\documentclass[10pt,a4paper]{article}

\usepackage[utf8]{inputenc}
\usepackage[T1]{fontenc}
\usepackage{lmodern}
\usepackage[margin=2.5cm]{geometry}
\usepackage{enumitem}
\usepackage{titlesec}
\usepackage{xcolor}
\usepackage[hyphens]{url}
\usepackage{xurl}
\usepackage[colorlinks,breaklinks=true]{hyperref}

% Link styling
\hypersetup{
    urlcolor=blue!70!black,
    linkcolor=blue!70!black,
    citecolor=blue!70!black,
}

% Section formatting
\titleformat{\section}{\large\bfseries\scshape}{}{0em}{}[\titlerule]
\titlespacing*{\section}{0pt}{1.5ex plus .5ex minus .2ex}{1ex plus .2ex}

\titleformat{\subsection}{\normalsize\bfseries}{}{0em}{}
\titlespacing*{\subsection}{0pt}{1ex plus .3ex minus .1ex}{0.5ex plus .1ex}

% Compact lists
\setlist[enumerate]{leftmargin=2em, itemsep=0.3ex, parsep=0ex, topsep=0.5ex}
\setlist[itemize]{leftmargin=2em, itemsep=0.3ex, parsep=0ex, topsep=0.5ex}

% Allow sloppy line breaking to prevent overfull hboxes with URLs
\sloppy

\begin{document}

% Header
\begin{center}
    {\LARGE\bfseries Gustav Nilsonne}\\[0.5em]
    Karolinska Institutet\\
    \href{mailto:gustav.nilsonne@ki.se}{gustav.nilsonne@ki.se}\\
    ORCID: \href{https://orcid.org/0000-0001-5273-0150}{0000-0001-5273-0150}
    \\
    GitHub: \href{https://github.com/GNilsonne}{GNilsonne}
     $\cdot$ Google Scholar: \href{https://scholar.google.com/citations?user=mXzy9UIAAAAJ}{mXzy9UIAAAAJ}
    
\end{center}

\vspace{0.5em}

%% ================================================================
%% DEGREES
%% ================================================================

\section{Degrees and Qualifications}
\begin{itemize}[nolistsep]

    \item MD, Karolinska Institutet, 2009

    \item PhD, Experimental Pathology, Karolinska Institutet, 2009: Experimental therapeutics against malignant mesothelioma: investigations in vitro. Full text: KI Open Archive \href{https://openarchive.ki.se/xmlui/handle/10616/38731}{[link]}

    \item Associate Professor (Docent), Karolinska Institutet, 2021

\end{itemize}


%% ================================================================
%% EMPLOYMENT
%% ================================================================

\section{Employment and Appointments}
\begin{itemize}[nolistsep]

    \item Military service in the mounted band of the Royal Lifeguard Dragoons, 2002

    \item Prosector, Department of Pathology, Karolinska University Hospital, 2005-2006

    \item Junior physician, geriatric and cognitive medicine, Karolinska University Hospital, 2010

    \item Postdoctoral Researcher, Karolinska Institutet, Department of Clinical Neuroscience, 2010-2011

    \item Visiting Postdoctoral Fellow, Donders Institute of Cognitive Neuroimaging, 2012

    \item Postdoctoral Researcher, Stockholm University, Stress Research Institute, 2012-2014

    \item Associated Postdoctoral Fellow, Karolinska Institutet, Department of Clinical Neuroscience, 2012-2014

    \item Researcher, Stockholm University, Stress Research Institute/Department of Psychology, 2014-

    \item Associated Researcher, Karolinska Institutet, Department of Clinical Neuroscience, 2014-2021

    \item Visiting Scholar, Stanford University, Department of Psychology, 2017

    \item Data Domain Specialist, Karolinska Institutet/Swedish National Data Service, 2017-2023

    \item National Coordinator of Data Domain Specialists, Swedish National Data Service, 2019-2023

    \item Visiting Fellow, QUEST Center, Berlin Institute of Health, Charité - Universitätsmedizin Berlin, 2020-2021

    \item Researcher, Karolinska Institutet, Department of Clinical Neuroscience, 2022-

    \item Visiting Scientist and Affiliate, Stanford University, Meta-Research Innovation Center at Stanford (METRICS), 2023-

    \item Senior Advisor, Swedish National Data Service, 2024-

\end{itemize}


%% ================================================================
%% PHD SUPERVISION
%% ================================================================

\section{Supervision of PhD Students}
\begin{itemize}[nolistsep]

    \item Michael Ingre, Stockholm University 2017, co-supervisor. Thesis: P-hacking in academic research: a critical review of the job strain model and of the association between night work and breast cancer in women. web: DiVA. \href{http://urn.kb.se/resolve?urn=urn:nbn:se:su:diva-141136}{[link]}

    \item Sandra Tamm, Karolinska Institutet 2019, co-supervisor. Thesis: A neuroimaging perspective on the emotional sleepy brain. web: openarchive.ki.se. \href{https://openarchive.ki.se/xmlui/handle/10616/46632}{[link]}

    \item Eva Hesselmark, Karolinska Institutet 2019, co-supervisor. Thesis: Pediatric acute neuropsychiatric syndrome : diagnosis, biomarkers and treatment. web: openarchive.ki.se. \href{https://openarchive.ki.se/xmlui/handle/10616/46741}{[link]}

    \item Hannibal Ölund Alonso, Karolinska Institutet, co-supervisor 2017-.

\end{itemize}


%% ================================================================
%% GRANTS
%% ================================================================

\section{Research Grants}
\begin{itemize}[nolistsep]

    \item Swedish Heart and Lung Foundation, 2005-2009, five grants for a total of 24 months of PhD student scholarships (grants no. 20051024, 20060003, 20070042, 20070209, and 20090797)

    \item Swedish Society for Medicine, 2011, 12 months postdoctoral scholarship

    \item Svensk Förening för Sömnforskning och Sömnmedicin, 2013, scholarship, 10 000 SEK

    \item Fredrik och Ingrid Thurings stiftelse, 2014, project grant, 125 000 SEK, grant no. 2014-00037

    \item Riksbankens Jubileumsfond, 2015, 3-year project grant, 4 053 000 SEK, grant no. P15-0310:1 web: swecris \href{https://www.vr.se/english/swecris.html#/project/P15-0310:1_RJ}{[link]}

    \item Fredrik och Ingrid Thurings stiftelse, 2015, project grant, 50 000 SEK, grant no. 2015-00170

    \item Swedish Medical Association, 2016, project grant, 50 000 SEK

    \item Rut och Arvid Wolff memorial foundation 2019, 121 500 SEK, grant no. 2019-01734

    \item Karolinska Institutet Research Foundation Grant, 2020, 86 000 SEK, grant no. 2020-01412

    \item Rut och Arvid Wolff memorial foundation, 2020, 121 500 SEK, grant no. 2020-02419

    \item Riksbankens Jubileumsfond, 2021, 3-year project grant, 4 159 000 SEK, grant no. P21-0384 web: swecris \href{https://www.vr.se/english/swecris.html#/project/P21-0384_RJ}{[link]}

    \item Region Stockholm, 2021, 3-year project grant, 2 000 000 SEK, grant no. FoUI-967372

    \item Horizon Europe, 2022, affiliated partner in project "Skills4EOSC", 3-year project grant, 4.5 person-months. web: CORDIS \href{https://cordis.europa.eu/project/id/101058527}{[link]}

    \item FTX Future Fund, 2022, 50 000 USD, committed but never disbursed

    \item Horizon Europe Marie Skłodowska-Curie Action "NAPS", host/supervisor to 2-year postdoc grant. web: CORDIS \href{https://cordis.europa.eu/project/id/101111182}{[link]}

    \item Horizon Europe, 2023, partner in project "Improving Reproducibility in Science (iRISE)", 3-year project grant, 21 person-months. web: CORDIS \href{https://cordis.europa.eu/project/id/101094853}{[link]}

    \item Horizon Europe Marie Skłodowska-Curie Action "ICTR", host/supervisor to 2-year postdoc grant. web: CORDIS \href{https://cordis.europa.eu/project/id/101152904}{[link]}

    \item Horizon Europe, 2024, partner in project "IP4OS Unpacking the possibilities of Intellectual Properties for Open Science", 2-year project grant, 27 person-months. web: CORDIS \href{https://cordis.europa.eu/project/id/101188026}{[link]}

    \item Svenska sällskapet för medicinsk forskning, 2025, host/supervisor to 3-year postdoc grant. web: SSMF \href{https://www.ssmf.se/forskare/oreste-affatato/}{[link]}

    \item Vinnova, 2025, 1-year project grant, 994 152 SEK. swecris \href{https://www.vr.se/english/swecris.html#/project/2025-00509_Vinnova}{[link]}

\end{itemize}


%% ================================================================
%% TEACHING
%% ================================================================

\section{Teaching}
\begin{itemize}[nolistsep]

    \item Teaching in the form of lectures, seminars, lab exercises, clinical supervision, and written and oral examination in the study programmes for medicine, dentistry, biomedical science, nursing, psychology, medical informatics, and nutrition, all at Karolinska Institutet, in the study programme for molecular cell biology, Södertörn University College, in the study programme for psychology, Uppsala University, and in the master's programme in Public Health at Stockholm University.

    \item Course assistant (amanuens):

    \item Pathology for dentistry students, Karolinska Institutet 2005

    \item Pathology for medical students, Karolinska Institutet 2005

    \item Cognitive neuroscience with applications in psychiatry, Karolinska Institutet 2011-2012

    \item Course leader:

    \item Introduction to brain imaging, PhD course, Karolinska Institutet 2014

    \item Cognitive neuroscience, elective course for medical students, Karolinska Institutet 2014-2015

    \item Ethics, Open science, and Reproducible research, PhD course, Stockholm University 2016-2018. Course materials: osf.io/ct9ad/

    \item Open science and Reproducible research, PhD course, Karolinska Institutet 2017-2023. Web: ki.se. Course materials: osf.io/gxbqs/

    \item Scientific Illustration and Visualization, PhD course, Karolinska Institutet 2023-

    \item Director of postgraduate studies (studierektor), Stress Research Institute, 2014-2019

    \item Supervisor to student thesis work:

    \item Semra Öz, Bachelor thesis in biomedical science, Karolinska Institutet, 2006: Undersökningar av experimentella cytostatika mot malignt mesoteliom (co-supervisor)

    \item Adam Szulkin, Summer research school in biomedical science, Karolinska Institutet, 2006: Phenotypic differences of Malignant Mesothelioma cells in sensitivity to Gemcitabine and Alimta

    \item Filip Mundt, Summer research school in biomedical science, Karolinska Institutet, 2006: Investigations of the Cytotoxicity of Carboplatin against Malignant Mesothelioma Cells and the Effects in Combination with Selenite

    \item Sandra Tamm, Degree project in medicine, Karolinska Institutet, 2010: Farmakologisk modulering av empati, ett beteendeförsök om effekterna av oxazepam – ett anxiolytiskt bensodiazepinpreparat – på empatiska responser

    \item Philip Little, Degree project in medicine, Karolinska Institutet, 2011: Effects of Midazolam and Flumazenil on Empathic Response – an Exploratory Study

    \item Madeleine Hedlund, Elective essay in psychology, 2015: The relationship between IL-6 and sleep – a meta-analysis

    \item Sahand Sharifzadeh, Degree project in medicine, Karolinska Institutet, 2016: The influence of maternal depression during pregnancy on neonatal fiber tract integrity url: openarchive.ki.se/xmlui/handle/10616/45432

    \item Erika Unsbo, Degree project in medicine, Karolinska Institutet, 2017: Motor abnormalities among children with PANS and/or PANDAS. doi: 10.6084/m9.figshare.4584034.v2

    \item Malin Blomberg, Degree project in medicine, Karolinska Institutet, 2017: Visuospatial ability in children with PANDAS and PANS. doi: 10.6084/m9.figshare.4591657.v1

    \item Arian Jafari and Bo Jenner, Degree project in psychology, Karolinska Institutet, 2017: Ethical review board-approved protocols and intent to use open practices in research on human subjects in Sweden. osf: osf.io/vgbw2

    \item Alia Lellouchi, Degree project in medicine, Karolinska Institutet, 2017: Cognitive and motor function in children with suspected PANS/PANDAS.

    \item Daniel Samsami, Degree project in medicine, Karolinska Institutet, 2017: Activation of ventral striatum during anticipation of reward – A pilot study using functional magnetic resonance imaging for a future experiment about psychopathy. doi: 10.6084/m9.figshare.5005238.v2

    \item Jakob Berefelt, Degree project in medicine, Karolinska Institutet, 2017: Activation of Anterior Insular Cortex and Anterior Medial Cingulate Cortex Associated with Direct Pain Perception and Empathy for Pain. doi: 10.6084/m9.figshare.5084557.v1

    \item Fredrik Hannfors, Degree project in medicine, Karolinska Institutet, 2017: Associations between Empathy for Pain, Emotional Mimicry and Emotional Contagion - Indicative Evidence from fMRI and EMG Paradigms. doi: 10.6084/m9.figshare.5307361.v1

    \item Jolinn Duvinger, Degree project in medicine, Karolinska Institutet, 2017: Effekter av COMT- polymorfism på empati för smärta - En studie med fMRI-data.

    \item Olea Schau Rybäck, Degree project in medicine, Karolinska Institutet, 2018: Reward anticipation and psychopathic traits - Association between Triarchic Psychopathy Measure subscales and activity in ventral striatum during Monetary Incentive Delay Task.

    \item Rickard Steijer, Degree project in medicine, Karolinska Institutet, 2018: Emotion regulation in self-reported disinhibition, boldness and meanness. doi: 10.17605/osf.io/8P9VG

    \item Johnny Soares, Degree project in medicine, Karolinska Institutet, 2019: The effects of sleep deprivation on vigilance and neural activation during the Sustained Attention to Response Task. doi: 10.31237/osf.io/e8bmt

    \item Poyan Karimi, Degree project in medicine, Karolinska Institutet, 2019: A functional magnetic resonance imaging study on sleep deprivation and its effect on empathy for pain and emotional regulation. doi: 10.31237/osf.io/r89q6

    \item Vivianne Jakobsson, Degree project in medicine, Karolinska Institutet, 2019: Resting state connectivity during total sleep deprivation: A functional magnetic resonance study doi: 10.31237/osf.io/3dtxz

    \item Sanna Graniittiaho, Degree project in medicine, Karolinska Institutet, 2020: Effects of sleep deprivation on empathy for pain and emotion regulation: A behavioral and neuroimaging study. doi: 10.31237/osf.io/z4gfy

    \item Lasse Anttila, Degree project in medicine, Karolinska Institutet, 2020: A Functional Magnetic Resonance Imaging Study on Global Signal Variability After One Night of Total sleep Deprivation doi: 10.31237/osf.io/dhqj2

    \item Anton Glasberg, Degree project in medicine, Karolinska Institutet, 2023: The role of Interleukin-6 as a biomarker in ADHD. doi: 10.31237/osf.io/4gfmk

    \item Dania Alkattan, Degree project in medicine, Karolinska Institutet, 2023

    \item Niclas Dittasiri, Degree project in medicine, Karolinska Institutet, 2023

    \item Sohaib Adem Osman, Degree project in medicine, Karolinska Institutet, 2023

    \item Iraida Suditu, Degree project in medicine, Karolinska Institutet, 2023

    \item Ash Mechlaoui, Degree project in medicine, Karolinska Institutet, 2024

    \item Sofie Emilson, Degree project in medicine, Karolinska Institutet, 2024

\end{itemize}


%% ================================================================
%% AWARDS
%% ================================================================

\section{Awards}
\begin{itemize}[nolistsep]

    \item Open Data Champion, SPARC Europe, 2017. web: openscholarchampions.eu/opendata/champion/dataforevidence/ \href{https://openscholarchampions.eu/opendata/champion/dataforevidence/}{[link]}

    \item Natur \& Kulturs populärvetenskapliga arbetsstipendium, 2018. web: nok.se \href{https://www.nok.se/priser-och-stipendier/stipendier/popularvetenskapliga-arbetsstipendier/natur--kulturs-popularvetenskapliga-pris-2018/}{[link]}

    \item British Neuroscience Association Credibility in Neuroscience Prize, 2023, awarded to the EEGManyLabs project. web: www.bna.org.uk/mediacentre/news/credibility-prize-2023/ \href{https://www.bna.org.uk/mediacentre/news/credibility-prize-2023/}{[link]}

\end{itemize}


%% ================================================================
%% PHD COMMITTEE
%% ================================================================

\section{PhD Thesis Committee Member / Faculty Opponent}
\begin{itemize}[nolistsep]

    \item University of Copenhagen, Tarang Sharma: Effects of selective serotonin reuptake inhibitors (SSRIs) and serotonin-norepinephrine reuptake inhibitors (SNRIs) on suicidality, violence, and quality of life. 2018-04-23.

\end{itemize}


%% ================================================================
%% ACADEMIC COMMISSIONS
%% ================================================================

\section{Academic Commissions of Trust}
\begin{itemize}[nolistsep]

    \item Member of committee for research and education (FoUU-kommittén), Karolinska University Hospital, 2005 (Student representative)

    \item Member of steering group for Karolinska Institutet Cancer network, 2004-2005 (Student representative)

    \item Speech writer to Karolinska Institutet president Harriet Wallberg-Henriksson, 2005-2007, with Carl Savage

    \item Member of Biotermgruppen, (Swedish Committee for Terminology in the Biomedical Sciences) 2010-, Vice president 2011

    \item Organiser, IgNobel tour events in Stockholm, yearly 2012-

    \item Co-founder of Nätverket för evidensbaserad policy (Swedish Network for Evidence-Based Policy) 2015, President 2016-2017, 2019-

    \item Organizer of Imaging the Sleepy Brain symposium, Karolinska Institutet 2014, with Torbjörn Åkerstedt Program: ki.se Abstracts: idear-net

    \item Ambassador for the Center for Open Science

    \item Participated in the Mozilla Festival 2014, London, UK

    \item Chairperson of the Committee for Badges for open practices, Center for Open Science, 2014-2017

    \item Organiser, Lecture series in collaboration with the Swedish Network for Psychedelic Science, Karolinska Institutet 2016

    \item Member of programme committee, SUHF conference on open science 2017

    \item Co-organiser of Stockholm BrainHack 2017 with Malin Sandström and Anders Eklund

    \item Member of Steering Committee for Thesis Commons

    \item Member of programme committee, Engaging Researchers in Good Data Management - A one-day conference sharing fresh ideas and good practices for creating and promoting effective research data practices, Cambridge 2017.

    \item Member of committee for evaluation of open access mandates, National Library of Sweden, 2017-2018. Report (in Swedish): kb.se

    \item Member of committee for the Data Engagement Management Award, 2018.

    \item Member of Research Data Office project team, Karolinska Institutet, 2018-2023

    \item Member of steering group for International Neuroinformatics Coordinating Facility Swedish Node, 2018

    \item Co-organiser of Stockholm BrainHack 2018 with Malin Sandström

    \item Member of editorial advisory board, Curie (Journal of the Swedish Research Council), 2018-2019

    \item Member of Research Data Alliance COVID-19 Clinical Working Group

    \item Representative of Karolinska Institutet to the European Open Science Cloud Association, 2021

    \item Member of national reference group for open research data, Swedish Research Council, 2021-

    \item Member of committee for a strategy beyond transformative deals with scientific publishers, The Association of Swedish Higher Education Institutions, 2021-2023

    \item Member of committee for implementing a national open science road map, The Association of Swedish Higher Education Institutions, 2021-2022

    \item Co-chair of Task Force for Research Careers, Credit, and Recognition, European Open Science Cloud Association (EOSC-A), 2021-2023

    \item Member of committee for research assessment in hiring and promotion, The Association of Swedish Higher Education Institutions, 2022-2024

    \item Member of EOSC-Future \& Research Data Alliance Artificial Intelligence \& Data Visitation Working Group, 2022-2023

    \item Member of Karolinska Institutet Open Science Working Group, 2022-

    \item Reviewer for scientific journals, books, and conferences:

    \item International Journal of Cancer

    \item Cancer Epidemiology, Biomarkers, and Prevention*

    \item BMC Cancer

    \item British Journal of Cancer*

    \item Journal of Hematology and Oncology*

    \item Science Translational Medicine

    \item Brain, Behavior, and Immunity*

    \item PLOS One*

    \item Scandinavian journal of Psychology*

    \item F1000 Research* (open peer review reports listed below)

    \item Organization for Human Brain Mapping 2017

    \item American Psychologist*

    \item Sleep Medicine Reviews*

    \item Journal of Inflammation Research*

    \item Organization for Human Brain Mapping 2018

    \item Translational Psychiatry*

    \item SLEEP*

    \item Human Brain Mapping*

    \item Molecular Psychiatry*

    \item Meta-Psychology

    \item Systematic Reviews*

    \item Psychiatry Research: Neuroimaging*

    \item European Sleep Medicine Textbook (Second Edition)*

    \item Royal Society Open Science*

    \item Meta-Psychology

    \item The Review of Psychology *Registered in Web of Science, ID: 1550920

    \item Reviewer for scientific grants:

    \item Horizon 2020 Marie Skłodowska Curie 2017

    \item Horizon 2020 Marie Skłodowska Curie 2018

    \item Horizon 2020 Marie Skłodowska Curie 2019

    \item Horizon 2020 Marie Skłodowska Curie 2020

    \item SPOKES Wellcome Trust Funded Translational Partnership Fellowships 2020, Berlin Institute of Health

    \item Horizon 2020 Future and Emerging Technologies (FET) Open 2020

    \item Horizon Europe Marie Skłodowska Curie 2021

    \item Norrsken Mind Foundation 2022

    \item ACX grants 2023

    \item ALF Region Skåne 2024

    \item Reviewer for scientific prizes

    \item British Neuroscience Association Credibility in Neuroscience Prize, 2024

    \item Reviewer of nonfiction books and book proposals

    \item Fri Tanke Förlag

    \item Elsevier

    \item SAGE Publications

    \item Guest editor of special issue of the Swedish Medical Journal on metascience, "Medicinsk forskning i fågelperspektiv", with Cathrine Axfors and Emmanuel Zavalis, 2023. web: lakartidningen.se/tag/tema-medicinsk-forskning-i-fagelperspektiv/

\end{itemize}


%% ================================================================
%% OTHER COMMISSIONS
%% ================================================================

\section{Other Commissions of Trust}
\begin{itemize}[nolistsep]

    \item Association of Swedish School Student Unions (Elevorganisationen i Sverige, now Sveriges Elevkårer)

    \item Elevorganisationen Stockholms Län, board member and treasurer, 1999-2000

    \item Elevorganisationen i Sverige, board member and coordinator for international cooperation, 2000-2001

    \item Elevorganisationen Stockholms Län, auditor, 2001-2002

    \item Organizing Bureau of European School Student Unions (OBESSU)

    \item Member of conference steering group "Languages in School", 2001

    \item Member of Internal Auditors' Commission, 2001-2005

    \item Advisor to the Board, 2005-2007

    \item Medical Students' Association at Karolinska Institutet (Medicinska Föreningen)

    \item Medical Students' Section, secretary, 2003

    \item Business committee, member 2003; president 2004

    \item Council (Kårfullmäktige), member 2003-2004

    \item Board, member 2004

    \item Students' Association of the Swedish Medical Association (Medicine Studerandes Förbund, now Läkarförbundet Student)

    \item Stockholm branch, board member 2003-2006

    \item Stockholm branch, vice president 2005

    \item National board, member 2003-2005

    \item Swedish Medical Association collective bargaining delegation (Förhandlingsdelegationen), 2004-2006

    \item Högskolerestauranger AB, board member, 2005-2006

    \item Stiftelsen Stockholms Studentbostäder (SSSB), board member, 2006

    \item Swedish Association for Sexuality Education (RFSU)

    \item RFSU Stockholm, board member, 2010

    \item RFSU Stockholm, president, 2011

    \item RFSU Stockholm, auditor, 2013-2018

    \item RFSU Researchers' Network, founding member, 2013-2018

    \item Nya Moderaterna: Medlem i arbetsgrupp för reformer, 2016-2017

    \item Swedish network for Evidence-Based Policy, founding member 2015, president 2018-

    \item Advisory Board, DeSci Foundation

    \item Advisory Board, NEKO Health

\end{itemize}


\section*{Research Outputs}

%% ================================================================
%% PUBLICATIONS
%% ================================================================

\section{Publications in Academic Journals}
\begin{enumerate}[nolistsep, label={[\arabic*]}]

    \item Nilsonne G, Sun X, Nyström C, Rundlöf A-K, Fernandes A, Björnstedt M, Dobra K. Selenite induces apoptosis in sarcomatoid malignant mesothelioma cells through oxidative stress. \textit{Free Radical Biology and Medicine} 2006. doi:~\href{https://doi.org/10.1016/j.freeradbiomed.2006.04.031}{10.1016/j.freeradbiomed.2006.04.031} [\href{http://hdl.handle.net/10616/47514}{preprint}]

    \item Rundlöf A-K, Fernandes A, Selenius M, Shariatgorji M, Nilsonne G, Ilag L, Dobra K and Björnstedt M. Quantification of alternative mRNA species and identification of thioredoxin reductase 1 isoforms in human tumor cells. \textit{Differentiation} 2007. doi:~\href{https://doi.org/10.1111/j.1432-0436.2006.00121.x}{10.1111/j.1432-0436.2006.00121.x}

    \item Nilsonne G, Olm E, Szulkin A, Mundt F, Stein A, Kosic B, Rundlöf A-K, Fernandes A P, Björnstedt M, Hjerpe A, and Dobra K. Phenotype-dependent apoptosis signalling in mesothelioma cells after selenite exposure. \textit{Journal of Experimental \& Clinical Cancer Research} 2009. doi:~\href{https://doi.org/10.1186/1756-9966-28-92}{10.1186/1756-9966-28-92}

    \item Zong F, Fthenou E, Wolmer N, Hollósi P, Kovalszky I, Szilák L, Mogler C, Nilsonne G, Tzanakakis G, and Dobra K. Syndecan-1 and FGF-2, but Not FGF Receptor-1, Share a Common Transport Route and Co-Localize with Heparanase in the Nuclei of Mesenchymal Tumor Cells. \textit{PLOS One} 2009. doi:~\href{https://doi.org/10.1371/journal.pone.0007346}{10.1371/journal.pone.0007346}

    \item Enqvist M, Nilsonne G, Hammarfjord O, Wallin R, Björkström N, Björnstedt M, Hjerpe A, Ljunggren H-G, Dobra K, Malmberg K-J, Carlsten, M. Selenite Induces Posttranscriptional Blockade of HLA-E Expression and Sensitizes Tumor Cells to CD94/NKG2A-Positive NK cells. \textit{Journal of Immunology} 2011. doi:~\href{https://doi.org/10.4049/jimmunol.1100610}{10.4049/jimmunol.1100610}

    \item Nilsonne G, Appelgren A, Axelsson J, Fredrikson M, Lekander M. Learning in a simple biological system: a pilot study of classical conditioning of human macrophages in vitro. \textit{Behavioral and brain functions} 2011. doi:~\href{https://doi.org/10.1186/1744-9081-7-47}{10.1186/1744-9081-7-47}

    \item Szulkin A, Nilsonne G, Mundt F, Wasik A, Souri P, Hjerpe A, Dobra K. Variation in Drug Sensitivity of Malignant Mesothelioma Cell Lines with Substantial Effects of Selenite and Bortezomib, Highlights Need for Individualized Therapy. \textit{PLOS One} 2013. doi:~\href{https://doi.org/10.1371/journal.pone.0065903}{10.1371/journal.pone.0065903}

    \item Darai-Ramqvist E, Nilsonne G, Flores-Staino C, Hjerpe A, Dobra K. Microenvironment-dependent phenotypic changes in a SCID mouse model for malignant mesothelioma. \textit{Frontiers in Oncology} 2013. doi:~\href{https://doi.org/10.3389/fonc.2013.00203}{10.3389/fonc.2013.00203} [\href{http://www.ncbi.nlm.nih.gov/geo/query/acc.cgi?acc=GSE48019}{data}]

    \item Mundt F*, Nilsonne G*, Arslan S, Csürös K, Hillerdal G, Yildirim H, Metintas M, Dobra K, Hjerpe A. Hyaluronan and N-ERC/Mesothelin as Key Biomarkers in a Specific Two-Step Model to Predict Pleural Malignant Mesothelioma. \textit{PLOS One} 2013. doi:~\href{https://doi.org/10.1371/journal.pone.0072030}{10.1371/journal.pone.0072030} [\href{http://dx.doi.org/10.5061/dryad.h122b}{data}]

    \item Mundt F, Heidari-Hamedari G, Nilsonne G, Metintas M, Hjerpe A, Dobra K. Diagnostic and Prognostic Value of Soluble Syndecan-1 for Pleural Malignancies. \textit{BioMed Research International} 2014. doi:~\href{https://doi.org/10.1155/2014/419853}{10.1155/2014/419853} [\href{http://dx.doi.org/10.5061/dryad.c42t7}{data}]

    \item Open Science Collaboration. Open Science Collaboration. \textit{Science} 2015. doi:~\href{https://doi.org/10.1126/science.aac4716}{10.1126/science.aac4716} [\href{https://osf.io/ezcuj}{data}]

    \item Nilsonne G, Tamm S, Månsson K, Åkerstedt T, Lekander M. Leukocyte telomere length and hippocampus volume: a meta-analysis. \textit{F1000Research} 2015. doi:~\href{https://doi.org/10.12688/f1000research.7198.1}{10.12688/f1000research.7198.1} [\href{http://dx.doi.org/10.5281/zenodo.31771}{data}]

    \item Lekander M, Karshikoff B, Johansson E, Soop A, Fransson P, Lundström J. N, Andreasson A, Ingvar M, Petrovic P, Axelsson J, Nilsonne G. Intrinsic functional connectivity of insular cortex and symptoms of sickness during acute experimental inflammation. \textit{Brain, Behavior, and Immunity} 2016. doi:~\href{https://doi.org/10.1016/j.bbi.2015.12.018}{10.1016/j.bbi.2015.12.018} [\href{http://dx.doi.org/10.5281/zenodo.35770}{data}]

    \item Anderson CJ, Bahník Š, Barnett-Cowan M, Bosco FA, Chandler J, Chartier CR, Cheung F, Christopherson CD, Cordes A, Cremata EJ, Della Penna N, Estel V, Fedor A, Fitneva SA, Frank MC, Grange JA, Hartshorne JK, Hasselman F, Henninger F, Jonas KJ, Lai CK, Levitan CA, Miller JK, Moore KS, Meixner JM, Munafò MR, Neijenhuijs KI, Nilsonne G, Nosek BA, Plessow F, Prenoveau JM, Ricker AA, Schmidt K, Spies JR, Stieger S, Strohminger N, Sullivan GB, van Aert RCM, van Assen MALM, van der Hulst M, Vanpaemel W, Vianello M, Voracek M, Zuni K. Response to a comment on “Estimating the reproducibility of psychological science”. \textit{Science} 2016. doi:~\href{https://doi.org/10.1126/science.aad9163}{10.1126/science.aad9163}

    \item Nilsonne G, Hilgard J, Lekander M, Arnberg F. Post-traumatic stress disorder and interleukin 6.  2016. doi:~\href{https://doi.org/10.1016/S2215-0366(16)00022-5}{10.1016/S2215-0366(16)00022-5} [\href{http://hdl.handle.net/10616/44992}{preprint} | \href{http://dx.doi.org/10.5281/zenodo.44726}{data}]

    \item Nilsonne G, Renberg A, Tamm S, Lekander M. Health at the ballot box: Disease threat does not predict attractiveness preference in British politicians. \textit{Royal Society Open Science} 2016. doi:~\href{https://doi.org/10.1098/rsos.160049}{10.1098/rsos.160049} [\href{http://dx.doi.org/10.5281/zenodo.12984}{data}]

    \item Sörman K, Nilsonne G, Howner K, Tamm S, Caman S, Wang H, Edens J, Gustavsson P, Lilienfeld S, Ingvar M, Petrovic P, Fischer H, Kristiansson M. Reliability and construct validity of the Psychopathic Personality Inventory-Revised in a Swedish non-criminal sample – a multimethod approach including psychophysiological correlates of empathy for pain. \textit{PLOS One} 2016. doi:~\href{https://doi.org/10.1371/journal.pone.0156570}{10.1371/journal.pone.0156570} [\href{http://dx.doi.org/10.5061/dryad.qh8c9}{data}]

    \item Nilsonne G, Lekander M, Åkerstedt T, Axelsson J, Ingre M. Diurnal variation of circulating interleukin-6: a meta-analysis. \textit{PLOS ONE} 2016. doi:~\href{https://doi.org/10.1371/journal.pone.0165799}{10.1371/journal.pone.0165799} [\href{http://dx.doi.org/10.1101/042507}{preprint} | \href{http://dx.doi.org/10.5281/zenodo.163325}{data}]

    \item Nilsonne G, Tamm S, Golkar A, Olsson A, Martin I, Petrovic P. Effects of 25 mg oxazepam on emotional mimicry and empathy for pain: a randomized controlled experiment. \textit{Royal Society Open Science} 2017. doi:~\href{https://doi.org/10.1098/rsos.160607}{10.1098/rsos.160607} [\href{http://dx.doi.org/10.6084/m9.figshare.1558201}{preprint} | \href{http://dx.doi.org/10.5281/zenodo.60385}{data}]

    \item Nilsonne G, Lekander M. Circulating interleukin-6 in Parkinson Disease. \textit{JAMA Neurology} 2017. doi:~\href{https://doi.org/10.1001/jamaneurol.2017.0037}{10.1001/jamaneurol.2017.0037} [\href{http://dx.doi.org/10.5281/zenodo.182505}{data}]

    \item Bejerot S, Nilsonne G, Humble M. Subcortical brain volume differences in participants with attention deficit hyperactivity disorder in children and adults.  2017. doi:~\href{https://doi.org/10.1016/S2215-0366(17)30160-8}{10.1016/S2215-0366(17)30160-8} [\href{http://dx.doi.org/10.5281/zenodo.345161}{data}]

    \item Nilsonne G, Tamm S, Schwarz J, Almeida R, Fischer H, Kecklund G, Lekander M, Fransson P, Åkerstedt T. Intrinsic brain connectivity after partial sleep deprivation in young and older adults: results from the Stockholm Sleepy Brain Study.  2017. doi:~\href{https://doi.org/10.1038/s41598-017-09744-7}{10.1038/s41598-017-09744-7} [\href{https://clinicaltrials.gov/ct2/show/NCT02000076}{prereg} | \href{http://dx.doi.org/10.1101/073494}{preprint} | \href{https://openneuro.org/dataset/ds000201/}{data}]

    \item Tamm S, Nilsonne G, Lamm C, Kecklund G, Petrovic P, Fischer H, Åkerstedt T, Lekander M. The effect of sleep restriction on empathy for pain: An fMRI study in younger and older adults.  2017. doi:~\href{https://doi.org/10.1038/s41598-017-12098-9}{10.1038/s41598-017-12098-9} [\href{https://clinicaltrials.gov/ct2/show/NCT02000076}{prereg} | \href{https://openneuro.org/dataset/ds000201/}{data}]

    \item Åkerstedt T, Lekander M, Nilsonne G, Tamm S, d'Onofrio P, Kecklund G, Fischer H, Schwarz J. Effects of late night short sleep on in-home polysomnography – relation to adult age and sex. \textit{Journal of Sleep Research} 2017. doi:~\href{https://doi.org/10.1111/jsr.12626}{10.1111/jsr.12626}

    \item Lakens D, et al. Justify Your Alpha: A Response to “Redefine Statistical Significance”. \textit{Nature Human Behaviour} 2018. doi:~\href{https://doi.org/10.1038/s41562-018-0311-x}{10.1038/s41562-018-0311-x} [\href{https://psyarxiv.com/9s3y6}{preprint} | \href{https://osf.io/by2kc/}{data}]

    \item Klein O, Hardwicke TE, Aust F, Breuer J, Danielsson H, Hofelich Mohr A, IJzerman H, Nilsonne G, Vanpaemel W, Frank MC. A Practical Guide for Transparency in Psychological Science. \textit{Collabra: Psychology}. doi:~\href{https://doi.org/10.1525/collabra.158}{10.1525/collabra.158} [\href{https://doi.org/10.17605/OSF.IO/RTYGM}{preprint}]

    \item Hardwicke TE, Mathur M, MacDonald K, Nilsonne G, Banks GC, Kidwell MC, Hofelich Mohr A, Clayton E, Yoon EJ, Tessler MH, Lenne RL, Altman S, Long B, Frank MC. Data availability, reusability, and analytic reproducibility: Evaluating the impact of a mandatory open data policy at the journal Cognition. \textit{Royal Society Open Science}. doi:~\href{https://doi.org/10.1098/rsos.180448}{10.1098/rsos.180448} [\href{https://doi.org/10.17605/OSF.IO/39CFB}{preprint}]

    \item Ingre M, Nilsonne G. Estimating statistical power, posterior probability and publication bias of psychological research using the observed replication rate. \textit{Royal Society Open Science}. doi:~\href{https://doi.org/10.1098/rsos.181190}{10.1098/rsos.181190} [\href{https://doi.org/10.6084/m9.figshare.5573113}{preprint} | \href{https://github.com/micing/publication_bias_psychology}{data}]

    \item Radun I, Nilsonne G, Radun J, Helgesson G, Kecklund G. Company employees as experimental participants in traffic safety research: prevalence and implications. \textit{Transportation Research Part F: Psychology and Behaviour} 2018. doi:~\href{https://doi.org/10.1016/j.trf.2018.10.008}{10.1016/j.trf.2018.10.008} [\href{http://hdl.handle.net/10616/46553}{preprint}]

    \item McGrath C, Nilsonne G. Data sharing in qualitative research: opportunities and concerns. \textit{MedEdPublish} 2018. doi:~\href{https://doi.org/10.15694/mep.2018.0000255.1}{10.15694/mep.2018.0000255.1} [\href{https://doi.org/10.31219/osf.io/ewyjb}{preprint}]

    \item Tamm S, Nilsonne G, Schwarz J, Golkar A, Kecklund G, Petrovic P, Fischer H, Åkerstedt T, Lekander M. Sleep restriction caused impaired emotional regulation without detectable brain activation changes - a functional magnetic resonance imaging study. \textit{Royal Society Open Science} 2019. doi:~\href{https://doi.org/10.1098/rsos.181704}{10.1098/rsos.181704} [\href{http://doi.org/10.1101/436048}{preprint}]

    \item Lodin K, Lekander M, Petrovic P, Nilsonne G, Hedman-Lagerlöf E, Andreasson A. Cross-sectional associations between inflammation, sickness behaviour, health anxiety and self-rated health in a Swedish primary care population. \textit{European Journal of Inflammation} 2019. doi:~\href{https://doi.org/10.1177/2058739219844357}{10.1177/2058739219844357}

    \item Landy J, et al. Crowdsourcing hypothesis tests: Making transparent how design choices shape research results. . doi:~\href{https://doi.org/10.1037/bul0000220}{10.1037/bul0000220} [\href{http://repository.essex.ac.uk/id/eprint/25784}{preprint}]

    \item Månsson KNT, Lindqvist D, Yang L, Svanborg C, Isung J, Nilsonne G, Bergman Nordgren L, El Alaou Si, Hedman-Lagerlöf E, Kraepelien M, Högström J, Andersson G, Boraxbekk C-J, Fischer H, Lavebratt C, Wolkowitz O, Furmark T. Improvement in indices of cellular protection after psychological treatment for social anxiety disorder.  2019. doi:~\href{https://doi.org/10.1038/s41398-019-0668-2}{10.1038/s41398-019-0668-2}

    \item Åkerstedt T, Lekander M, Nilsonne G, Tamm S, d'Onofrio P, Kecklund G, Fischer H, Schwarz J, Petrovic P, Månsson KNT. Gray matter volume correlates of sleepiness: a voxel–based morphometry study in younger and older adults. \textit{Nature and Science of Sleep} 2020. doi:~\href{https://doi.org/10.2147/NSS.S240493}{10.2147/NSS.S240493}

    \item Evans T, Johannes N, Winska J, Glińska-Neweś A, van Stekelenburg A, Nilsonne G, Dean L, Fido D, Galloway G, Jones S, Masson I, Soares A, Steptoe-Warren G, Thompson N, d'Angelo Ungson N. Exploring the Consistency and Value of Humour Style Clusters. \textit{Comprehensive Results in Social Psychology} 2020. doi:~\href{https://doi.org/10.1080/23743603.2020.1756239}{10.1080/23743603.2020.1756239} [\href{https://doi.org/10.31234/osf.io/b8g9z}{preprint}]

    \item Botvinik-Nezer R, et al. Variability in the analysis of a single neuroimaging dataset by many teams. \textit{Nature} 2020. doi:~\href{https://doi.org/10.1038/s41586-020-2314-9}{10.1038/s41586-020-2314-9} [\href{https://doi.org/10.1101/843193}{preprint}]

    \item Koenig J, et al. Cortical Thickness and Resting State Cardiac Function Across the Lifespan: A Cross-Sectional Pooled Mega Analysis. \textit{Psychophysiology} 2021. doi:~\href{https://doi.org/10.1111/psyp.13688}{10.1111/psyp.13688} [\href{https://doi.org/10.31219/osf.io/btjpw}{preprint}]

    \item Tierney W, Hardy J, Ebersole C, Viganola D, Clemente E, Gordon M, Hoogeveen S, Haaf J, Dreber A, Johannesson M, Pfeiffer T, Huang J L, Vaugh L A, DeMarree K G, Igou E, Chapman H, Gantman A, Vanaman M, Wylie J, Storbeck J, Andreychik M R, McPhetres J Culture and Work Forecasting Collaboration, Uhlmann E L. Culture and Work Forecasting Collaboration. \textit{Journal of Experimental Social Psychology.} 2020. doi:~\href{https://doi.org/10.1016/j.obhdp.2020.07.002}{10.1016/j.obhdp.2020.07.002}

    \item Tamm S, Schwarz J, Thuné H, Kecklund G, Petrovic P, Åkerstedt T, Fischer H, Lekander M, Nilsonne G. A combined fMRI and EMG study of emotional contagion following partial sleep deprivation in young and older humans.  2020. doi:~\href{https://doi.org/10.1038/s41598-020-74489-9}{10.1038/s41598-020-74489-9} [\href{https://doi.org/10.31234/osf.io/k6acx}{preprint}]

    \item Nilsonne G, Harrell F. EEG-based model and antidepressant response.  2021. doi:~\href{https://doi.org/10.1038/s41587-020-00768-5}{10.1038/s41587-020-00768-5} [\href{https://doi.org/10.31234/osf.io/hqxzd}{preprint}]

    \item Sorjonen K, Nilsonne G, Falkstedt D, Hemmingson T, Ingre M. A comparison of models with weight, height, and BMI as predictors of mortality. \textit{Obesity Science and Practice} 2020. doi:~\href{https://doi.org/10.1002/osp4.473}{10.1002/osp4.473} [\href{https://doi.org/10.31234/osf.io/xamzc}{preprint}]

    \item Tierney W, Hardy JH, Ebersole CR, Leavitt K, Viganola D, Clemente EG, Gordon M, Dreber A, Johannesson M, Pfeiffer T, Hiring Decisions Forecasting Collaboration, Uhlmann EL. Hiring Decisions Forecasting Collaboration. \textit{Organizational Behavior and Human Decision Processes} 2020. doi:~\href{https://doi.org/10.1016/j.obhdp.2020.07.002}{10.1016/j.obhdp.2020.07.002}

    \item Koba C, Notaro G, Tamm S, Nilsonne G, Hasson U. Spontaneous eye-movements during eyes-open rest reduce resting-state-network modularity by increasing visual-sensorimotor connectivity.  2020. doi:~\href{https://doi.org/10.1162/netn\_a\_00186}{10.1162/netn\\_a\\_00186} [\href{https://doi.org/10.1101/2020.05.18.100669}{preprint}]

    \item Pavlov YG, Adamian N, Appelhoff S, Arvaneh M, Christopher Benwell PD, Beste C, Bland A, Bradford DE, Bublatzky F, Busch N, Clayson PE, Cruse D, Czeszumski A, Almenberg AD, Dumas G, Ehinger BV, Ganis G, He X, Hinojosa JA, Huber-Huber C, Inzlicht M, Jack BN, Johannesson M, Jones R, Kaltwasser L, Karimi-Rouzbahani H, Keil A, König P, Kouara L, Kulke L, Ladouceur C, Langer N, Liesefeld HR, Luque D, MacNamara A, Mudrik L, Muthuraman M, Neal L, Nilsonne G, Galán JGN, Ocklenburg S, Oostenveld R, Pernet DC, Pourtois G, Ruzzoli M, Sass S, Schaefer A, Senderecka M, Snyder JS, Tamnes CK, Tognoli E, Vugt MK van, Verona E, Vloeberghs R, Welke D, Wessel J, Zakharov I, Mushtaq F. \#EEGManyLabs: Investigating the Replicability of Influential EEG Experiments. \textit{Cortex} 2021. doi:~\href{https://doi.org/10.1016/j.cortex.2021.03.013}{10.1016/j.cortex.2021.03.013} [\href{https://doi.org/10.31234/osf.io/528nr}{preprint}]

    \item Nilsonne G, Tamm S, Golkar A, Olsson A, Sörman K, Howner K, Kristiansson M, Ingvar M, Petrovic P. Oxazepam and cognitive reappraisal: a randomised experiment. . doi:~\href{https://doi.org/10.1101/656355}{10.1101/656355} [\href{https://doi.org/10.1101/656355}{preprint} | \href{https://doi.org/10.5281/zenodo.3236076}{data}]

    \item Sorjonen K, Nilsonne G, Melin B, Ingre M. The new Accounting for Expected Adjusted Effect test (AEAE test) has higher positive predictive value than a zero-order significance test. \textit{BMC Research Notes}. doi:~\href{https://doi.org/10.1186/s13104-021-05545-4}{10.1186/s13104-021-05545-4} [\href{http://doi.org/10.31234/osf.io/fp8g3}{preprint}]

    \item Gau R, et al. Brainhack: Developing a culture of open, inclusive, community-driven neuroscience. \textit{Neuron} 2021. doi:~\href{https://doi.org/10.1016/j.neuron.2021.04.001}{10.1016/j.neuron.2021.04.001} [\href{http://doi.org/10.31234/osf.io/rytjq}{preprint}]

    \item Hallinan D, Bernier A, Cambon-Thomsen A, Crawley F P, Dimitrova D, Medeiros C, Nilsonne G, Parker S, Pickering B, Rennes S. International Transfers of Health Research Data Following Schrems II: A Problem in Need of a Solution. \textit{European Journal of Human Genetics} 2021. doi:~\href{https://doi.org/10.1038/s41431-021-00893-y}{10.1038/s41431-021-00893-y} [\href{http://doi.org/10.2139/ssrn.3688392}{preprint}]

    \item Schweinsberg M, et al. Same data, different conclusions: Radical dispersion in empirical results when independent analysts operationalize and test the same hypothesis. \textit{Organizational Behavior and Human Decision Processes} 2021. doi:~\href{https://doi.org/10.1016/j.obhdp.2021.02.003}{10.1016/j.obhdp.2021.02.003} [\href{http://repository.essex.ac.uk/id/eprint/29938}{preprint}]

    \item Tahmasian M, Aleman A, Andreassen OA, Arab Z, Baillet M, Benedetti F, Bresser T, Bright J, Chee MWL, Chylinski D, Cheng W, Deantoni M, Dresler M, Eickhoff SB, Eickhoff CR, Elvsåshagen T, Feng J, Foster‐Dingley JC, Ganjgahi H, Grabe HJ, Groenewold NA, Ho TC, Hong SB, Houenou J, Irungu B, Jahanshad N, Khazaie H, Kim H, Koshmanova E, Kocevska D, Kochunov P, Lakbila‐Kamal O, Leerssen J, Li M, Luik AI, Muto V, Narbutas J, Nilsonne G, O’Callaghan VS, Olsen A, Osorio RS, Poletti S, Poudel G, Reesen JE, Reneman L, Reyt M, Riemann D, Rosenzweig I, Rostampour M, Saberi A, Schiel J, Schmidt C, Schrantee A, Sciberras E, Silk TJ, Sim K, Smevik H, Soares JC, Spiegelhalder K, Stein DJ, Talwar P, Tamm S, Teresi G l, Valk SL, Someren EV, Vandewalle G, Egroo MV, Völzke H, Walter M, Wassing R, Weber FD, Weihs A, Westlye LT, Wright MJ, Wu M-J, Zak N, Zarei M. ENIGMA-Sleep: Challenges, opportunities, and the road map. \textit{Journal of Sleep Research}. doi:~\href{https://doi.org/gjtqz6}{gjtqz6}

    \item Wang K, et al. A multi-country test of brief reappraisal interventions on emotions during the COVID-19 pandemic. \textit{Nature Human Behavior} 2021. doi:~\href{https://doi.org/10.1038/s41562-021-01173-x}{10.1038/s41562-021-01173-x} [\href{https://doi.org/10.31234/osf.io/m4gpq}{preprint}]

    \item Austin CC, Bernier A, Bezuidenhout L, Bicarregui J, Biro T, Cambon-Thomsen A, Russo Carroll S, Cournia Z, Dabrowski PW, Diallo G, Duflot T, Garcia L, Gesing S, Gonzalez-Beltran A, Gururaj A, Harrower N, Lin D, Medeiros C, Méndez E, Meyers N, Mietchen D, Nagrani R, Nilsonne G, Parker S, Pickering B, Pienta A, Polydoratou P, Psomopoulos F, Rennes S, Rowe R, Sansone S-A, Shanahan H, Sitz L, Stocks J, Tovani-Palone MR, Uhlmansiek M. Fostering global data sharing: highlighting the recommendations of the Research Data Alliance COVID-19 working group. \textit{Wellcome Open Research} 2021. doi:~\href{https://doi.org/10.12688/wellcomeopenres.16378.2}{10.12688/wellcomeopenres.16378.2}

    \item Sorjonen K, Falkstedt D, Wallin AS, Melin B, Nilsonne G. Dangers of residual confounding: A cautionary tale featuring cognitive ability, socioeconomic background, and education. \textit{BMC Psychology} 2021. doi:~\href{https://doi.org/10.1186/s40359-021-00653-z}{10.1186/s40359-021-00653-z} [\href{https://doi.org/10.31234/osf.io/esfx5}{preprint}]

    \item Sorjonen K, Nilsonne G, Ingre M, Melin B. Spurious correlations in research on ability tilt. \textit{Personality and Individual Differences} 2022. doi:~\href{https://doi.org/10.1016/j.paid.2021.111268}{10.1016/j.paid.2021.111268} [\href{https://doi.org/10.31234/osf.io/cg8hd}{preprint}]

    \item Aczel B, Szaszi B, Nilsonne G, van den Akker OR, Albers CJ, van Assen MA, Bastiaansen JA, Benjamin D, Boehm U, Botvinik-Nezer R, Bringmann LF, Busch NA, Caruyer E, Cataldo AM, Cowan N, Delios A, van Dongen NN, Donkin C, van Doorn JB, Dreber A, Dutilh G, Egan GF, Gernsbacher MA, Hoekstra R, Hoffmann S, Holzmeister F, Huber J, Johannesson M, Jonas KJ, Kindel AT, Kirchler M, Kunkels YK, Lindsay DS, Mangin J-F, Matzke D, Munafò MR, Newell BR, Nosek BA, Poldrack RA, van Ravenzwaaij D, Rieskamp J, Salganik MJ, Sarafoglou A, Schonberg T, Schweinsberg M, Shanks D, Silberzahn R, Simons DJ, Spellman BA, St-Jean S, Starns JJ, Uhlmann EL, Wicherts J, Wagenmakers E-J. Consensus-based guidance for conducting and reporting multi-analyst studies. \textit{ELife} 2021. doi:~\href{https://doi.org/10.7554/eLife.72185}{10.7554/eLife.72185} [\href{https://doi.org/10.31222/osf.io/5ecnh}{preprint}]

    \item Sorjonen K, Melin B, Nilsonne G. Lord’s paradox in latent change score modeling: an example involving facilitating longitudinal effects between intelligence and academic achievement. \textit{Personality and Individual Differences} 2022. doi:~\href{https://doi.org/10.1016/j.paid.2021.111520}{10.1016/j.paid.2021.111520}

    \item Sorjonen K, Nilsonne G, Ingre M, Melin B. Curiosity might not help after all: Predicted trajectories for need for cognition and anxiety and depression symptoms based on findings by Zainal and Newman (2022). \textit{Journal of Affective Disorders} 2022. doi:~\href{https://doi.org/10.1016/j.jad.2022.01.109}{10.1016/j.jad.2022.01.109} [\href{https://doi.org/10.31234/osf.io/dh2zm}{preprint}]

    \item Delios A, Clemente EG, Wu T, Tan H, Wang Y, Gordon M, Viganola D, Chen Z, Dreber A, Johannesson M, Pfeiffer T, Generalizability Tests Forecasting Collaboration, Uhlmann E. Generalizability Tests Forecasting Collaboration. \textit{Proceedings of the National Academy of Sciences} 2022. doi:~\href{https://doi.org/10.1073/pnas.2120377119}{10.1073/pnas.2120377119}

    \item Helgesson G, Radun I, Radun J, Nilsonne G. Editors publishing in their own journals – a systematic review of prevalence and a discussion of normative aspects. . doi:~\href{https://doi.org/10.1002/leap.1449}{10.1002/leap.1449} [\href{https://doi.org/10.31222/osf.io/dc6je}{preprint}]

    \item Bourget M-H, Kamentsky L, Ghosh S, Mazzamuto G, Lazari A, Markiewicz CJ, Oostenveld R, Niso G, Halchenko YO, Lipp I, Takerkart S, Toussaint P-J, Khan AR, Nilsonne G, Castelli FM, BIDS Maintainers, Cohen-Adad J. Microscopy-BIDS: an extension to the Brain Imaging Data Structure for Microscopy Data. \textit{Frontiers in Neuroscience} 2022. doi:~\href{https://doi.org/10.3389/fnins.2022.871228}{10.3389/fnins.2022.871228}

    \item Psychological Science Accelerator Self-Determination Theory Collaboration. Psychological Science Accelerator Self-Determination Theory Collaboration. \textit{Proceedings of the National Academy of Sciences} 2022. doi:~\href{https://doi.org/10.1073/pnas.2111091119}{10.1073/pnas.2111091119} [\href{https://doi.org/10.31234/osf.io/n3dyf}{preprint}]

    \item Sorjonen K, Nilsonne G, Ingre M, Melin B. Regression to the mean in latent change score models: an example involving breastfeeding and intelligence. \textit{BMC Pediatrics} 2022. doi:~\href{https://doi.org/10.1186/s12887-022-03349-4}{10.1186/s12887-022-03349-4}

    \item Dorison C, et al. In COVID-19 health messaging, loss framing increases anxiety with little-to-no concomitant benefits: Experimental evidence from 84 countries. \textit{Affective science} 2022. doi:~\href{https://doi.org/10.1007/s42761-022-00128-3}{10.1007/s42761-022-00128-3}

    \item Hoogeveen S, et al. A many-analysts approach to the relation between religiosity and well-being. \textit{Religion, Brain \& Behavior} 2022. doi:~\href{https://doi.org/10.1080/2153599X.2022.2070255}{10.1080/2153599X.2022.2070255}

    \item TARG Meta-Research Group \& Collaborators. TARG Meta-Research Group \& Collaborators. \textit{Royal Society Open Science} 2014. doi:~\href{https://doi.org/10.1098/rsos.220142}{10.1098/rsos.220142} [\href{https://doi.org/10.1101/2022.01.18.22269507}{preprint}]

    \item Lindsäter E, Svärdman F, Wallert J, Ivanova EN, Söderholm A, Fondberg R, Nilsonne G, Cervenka S, Lekander M, Rück C. Exhaustion Disorder: A Scoping Review of Research on a Recently Introduced Stress-Related Diagnosis. \textit{BJPsych Open} 2022. doi:~\href{https://doi.org/10.1192/bjo.2022.559}{10.1192/bjo.2022.559} [\href{https://doi.org/10.31234/osf.io/m4w9x}{preprint}]

    \item Sorjonen K, Nilsonne G, Ingre M, Melin B. Questioning the vulnerability model: Prospective associations between low self-esteem and subsequent degree of depression may be spurious. \textit{Journal of Affective Disorders} 2022. doi:~\href{https://doi.org/10.1016/j.jad.2022.08.003}{10.1016/j.jad.2022.08.003} [\href{https://doi.org/10.31234/osf.io/4kcma}{preprint}]

    \item Drude NI, Martinez-Gamboa L, Danziger M, Collazo A, Kniffert S, Wiebach J, Nilsonne G, Konietschke F, Piper SK, Pawel S, Micheloud C, Held L, Frommlet F, Segelcke D, Pogatzki-Zahn EM, Voelkl B, Friede T, Brunner E, Dempfle A, Haller B, Jung MJ, Riecken LB, Kuhn H-G, Tenbusch M, Higuita LMS, Remarque EJ, Grüninger-Egli SL, Manske K, Kobold S, Rivalan M, Wedekind L, Wilcke JC, Boulesteix A-L, Meinhardt MW, Spanagel R, Hettmer S, von Lüttichau I, Regina C, Dirnagl U, Toelch U. Planning preclinical confirmatory multicenter trials to strengthen translation from basic to clinical research – a multi-stakeholder workshop report. \textit{Translational Medicine Communications} 2022. doi:~\href{https://doi.org/10.1186/s41231-022-00130-8}{10.1186/s41231-022-00130-8}

    \item Sorjonen K, Nilsonne G, Melin B, Ingre M. Uncertain inference in random intercept cross-lagged panel models: An example involving need for cognition and anxiety and depression symptoms. \textit{Personality and Individual Differences} 2023. doi:~\href{https://doi.org/10.1016/j.paid.2022.111925}{10.1016/j.paid.2022.111925} [\href{https://doi.org/10.31234/osf.io/2a5gp}{preprint}]

    \item Buchanan EM, Lewis SC, Paris B, Forscher PS, Pavlacic JM, Beshears JE, et al. The Psychological Science Accelerator’s COVID-19 rapid-response dataset.  2023. doi:~\href{https://doi.org/10.1038/s41597-022-01811-7}{10.1038/s41597-022-01811-7}

    \item Sorjonen K, Ingre M, Nilsonne G, Melin B. Dangers of including outcome at baseline as a covariate in latent change score models: Results from simulations and empirical re-analyses. \textit{Heliyon} 2023. doi:~\href{https://doi.org/10.1016/j.heliyon.2023.e15746}{10.1016/j.heliyon.2023.e15746} [\href{https://doi.org/10.31234/osf.io/2sxja}{preprint}]

    \item Sorjonen K, Ingre M, Nilsonne G,. Melin B. Further arguments that ability tilt correlations are spurious: A reply to Coyle (2022). \textit{Intelligence} 2022. doi:~\href{https://doi.org/10.1016/j.intell.2022.101706}{10.1016/j.intell.2022.101706} [\href{https://doi.org/10.31234/osf.io/7vhnz}{preprint}]

    \item Nilsonne G, Johannesson M, Dreber A. Olika slutsatser från samma data. \textit{Läkartidningen} 2023.

    \item Brembs B, Huneman P, Schönbrodt F, Nilsonne G, Susi T, Siems R, et al. Replacing academic journals. \textit{Royal Society Open Science} 2023. doi:~\href{https://doi.org/10.1098/rsos.230206}{10.1098/rsos.230206} [\href{https://doi.org/10.5281/zenodo.7974116}{preprint}]

    \item Sorjonen K, Ingre M, Melin B, Nilsonne G. Unmasking artifactual links: A reanalysis reveals No direct causal relationship between self-esteem and quality of social relations. \textit{Heliyon} 2023.

    \item Trübutschek D, Yang YF, Gianelli C, Cesnaite E, Fischer NL, Vinding MC, et al. EEGManyPipelines: A Large-scale, Grassroots Multi-analyst Study of Electroencephalography Analysis Practices in the Wild. \textit{Journal of Cognitive Neuroscience} 2024. doi:~\href{https://doi.org/10.1162/jocn\_a\_02087}{10.1162/jocn\\_a\\_02087} [\href{https://doi.org/10.31222/osf.io/jq342}{preprint}]

    \item Sorjonen K, Melin B, Nilsonne G. Spurious heritability of ability tilts: A comment on Coyle et al. (2023). \textit{Personality and Individual Differences} 2023. doi:~\href{https://doi.org/10.31234/osf.io/kds5c}{10.31234/osf.io/kds5c} [\href{https://doi.org/10.31234/osf.io/kds5c}{preprint}]

    \item Sorjonen K, Nilsonne G, Melin B. Distorted meta-analytic findings on peer influence: A reanalysis. \textit{Heliyon} 2023. doi:~\href{https://doi.org/110.1016/j.heliyon.2023.e21458}{110.1016/j.heliyon.2023.e21458} [\href{https://doi.org/10.31234/osf.io/szkt8}{preprint}]

    \item Vlasceanu M, et al. Addressing Climate Change with Behavioral Science: A Global Intervention Tournament in 63 Countries. . doi:~\href{https://doi.org/10.1126/sciadv.adj5778}{10.1126/sciadv.adj5778} [\href{https://doi.org/10.31234/osf.io/cr5at}{preprint}]

    \item Sorjonen K, Nilsonne G, Ingre M, Melin B. Breastfeeding, cognitive ability, and residual confounding: A comment on Pereyra-Elìas et al. (2022). \textit{PLoS ONE} 2022. doi:~\href{https://doi.org/10.1371/journal.pone.0297216}{10.1371/journal.pone.0297216} [\href{https://doi.org/10.31234/osf.io/gkftq}{preprint}]

    \item Poldrack RA, et al. The past, present, and future of the brain imaging data structure (BIDS).  2024. doi:~\href{https://doi.org/10.1162/imag\_a\_00103}{10.1162/imag\\_a\\_00103} [\href{https://doi.org/10.48550/arXiv.2309.05768}{preprint}]

    \item Weissgerber TL, Gazda MA, Nilsonne G, et al. Understanding the provenance and quality of methods is essential for responsible reuse of FAIR data.  2024.

    \item Sorjonen K, Melin B, Nilsonne G. Inconclusive evidence for an increasing effect of maternal supportiveness on childhood intelligence in Dunkel et al. (2023): A simulated reanalysis. \textit{Intelligence} 2023. doi:~\href{https://doi.org/10.1016/j.intell.2024.101815}{10.1016/j.intell.2024.101815} [\href{https://doi.org/10.31234/osf.io/945hb}{preprint}]

    \item Eriksson K, Sorjonen K, Falkstedt D, Melin B, Nilsonne G. A formal model accounting for measurement reliability shows attenuated effect of higher education on intelligence in longitudinal data. \textit{Royal Society Open Science} 2024. doi:~\href{https://doi.org/10.1098/rsos.230513}{10.1098/rsos.230513} [\href{https://doi.org/10.31234/osf.io/xajz8}{preprint}]

    \item Axfors C, Nilsonne G. Registrering och rapportering av kliniska studier kan bli bättre. \textit{Läkartidningen} 2024.

    \item Sarafoglou A, et al. Subjective Evidence Evaluation Survey For Multi-Analyst Studies. \textit{Royal Society Open Science} 2024. doi:~\href{https://doi.org/10.1098/rsos.240125}{10.1098/rsos.240125} [\href{https://doi.org/10.31234/osf.io/mxje3}{preprint}]

    \item Mushtaq F, et al. One hundred years of EEG for brain and behaviour research. \textit{Nature Human Behaviour} 2024. doi:~\href{https://doi.org/10.1038/s41562-024-01941-5}{10.1038/s41562-024-01941-5}

    \item Doell KC, et al. The International Climate Psychology Collaboration: Climate change-related data collected from 63 countries.  2024. doi:~\href{https://doi.org/10.1038/s41597-024-03865-1}{10.1038/s41597-024-03865-1}

    \item Lindberg O, et al. Altered Empathy Processing in Frontotemporal Dementia. \textit{JAMA Network Open} 2024. doi:~\href{https://doi.org/10.1101/2024.03.21.586051}{10.1101/2024.03.21.586051} [\href{https://doi.org/10.1001/jamanetworkopen.2024.48601}{preprint}]

    \item Shrike BT, Nilsonne G, Weisser V, Aimbetov R. Longevity biomarkers: applications and limitations. \textit{ResearchHub}. doi:~\href{https://doi.org/10.55277/ResearchHub.rb3ck05l}{10.55277/ResearchHub.rb3ck05l}

    \item Sennerstam V, et al. Exhaustion Disorder in Primary Care: A Comparison With Major Depressive Disorder and Adjustment Disorder. \textit{Scandinavian Journal of Psychology} 2024. doi:~\href{https://doi.org/10.1111/sjop.13087}{10.1111/sjop.13087}

    \item Nilsonne G, et al. Results reporting for clinical trials led by medical universities and university hospitals in the Nordic countries was often missing or delayed. \textit{Journal of Clinical Epidemiology} 2025. doi:~\href{https://doi.org/10.1016/j.jclinepi.2025.111710}{10.1016/j.jclinepi.2025.111710}

    \item Gould E, et al. Same data, different analysts: variation in effect sizes due to analytical decisions in ecology and evolutionary biology. \textit{BMC Biology.} 2025. doi:~\href{https://doi.org/10.1186/s12915-024-02101-x}{10.1186/s12915-024-02101-x} [\href{https://doi.org/10.32942/X2GG62}{preprint}]

\end{enumerate}


%% ================================================================
%% PREPRINTS
%% ================================================================

\section{Preprints}
\begin{enumerate}[nolistsep, label={[\arabic*]}]

    \item Nilsonne G, Tamm S, d’Onofrio P, Thuné HÅ, Schwarz J, Lavebratt C, Liu, JJ, Månsson KNT, Sundelin T, Axelsson J, Fransson P, Kecklund G, Fischer H, Lekander M, Åkerstedt Torbjörn. A multimodal brain imaging dataset on sleep deprivation in young and old humans. [\href{https://openneuro.org/dataset/ds000201/}{data}]

    \item Ingre M, Sorjonen K, Nilsonne G. A cocktail fallacy in research with composite variables: a commentary on job strain and effort–reward imbalance (ERI) theories with updated meta-analyses on health outcomes. doi:~\href{https://doi.org/10.31234/osf.io/9hyxj}{10.31234/osf.io/9hyxj}

    \item Tennant J, Agarwal R, Baždarić K, Brassard D, Crick T, Dunleavy DJ, Evans TR, Gardner N, Gonzalez-Marquez M, Graziotin D, Greshake Tzovaras B, Gunnarsson D, Havemann J, Hosseini M, Katz DS, Knöchelmann M, Madan CR, Manghi P, Marocchino A, Masuzzo P, Murray-Rust P, Narayanaswamy S, Nilsonne G, Pacheco-Mendoza J, Penders B, Pourret O, Rera M, Samuel J, Steiner T, Stojanovski J, Uribe-Tirado A, Vos R, Worthington S, Yarkoni T. A tale of two “opens”: Intersections between Free and Open Source Software and Open Scholarship.

    \item Nilsonne G, Schwarz J, Kecklund G, Petrovic P, Fischer H, Åkerstedt T, Lekander M, Tamm S. Empirical evidence for a three-level model of emotional contagion, empathy and emotional regulation. doi:~\href{https://doi.org/10.31234/osf.io/ake34}{10.31234/osf.io/ake34}

    \item van der Hoven M, Asaduzzaman M, Evans N, Ike C, Kalichman M, Kniffert S, et al. Seven challenges for research integrity education: current status and recommendations. doi:~\href{https://doi.org/10.31219/osf.io/5w9kg}{10.31219/osf.io/5w9kg} [\href{https://doi.org/10.31219/osf.io/5w9kg}{preprint}]

    \item Sorjonen K, Nilsonne G, Ghaderi A, Melin B. Spurious effects in random-intercept cross-lagged panel models: Results from simulations and reanalyses of data on self-esteem and problematic eating behaviors used by Beckers et al. (2023). doi:~\href{https://doi.org/10.31234/osf.io/hferw}{10.31234/osf.io/hferw}

    \item Nilsonne G, et al. Results reporting for clinical trials led by medical universities and university hospitals in the Nordic countries was often missing or delayed. \textit{medRxiv}. doi:~\href{https://doi.org/10.1101/2024.02.04.24301363}{10.1101/2024.02.04.24301363}

    \item Sorjonen K, Arshamian A, Lager E, Nilsonne G,Melin B. In aqua veritas: Triangulation of cross-lagged effects for improved causal inference. doi:~\href{https://doi.org/10.31234/osf.io/ztwy6}{10.31234/osf.io/ztwy6}

    \item Sorjonen K, Arshamian A, Lager E, Nilsonne G,Melin B. A taxonomy of treatment effects in data with two waves of measurement and a promotion of triangulation. doi:~\href{https://doi.org/10.31234/osf.io/q47s6}{10.31234/osf.io/q47s6}

    \item Custer RM, Lynch KM, Barisano G, Herting MM, Åkerstedt T, Nilsonne G, et al. Effects of one-night partial sleep deprivation on perivascular space volume fraction: Findings from the Stockholm Sleepy Brain Study. doi:~\href{https://doi.org/10.1101/2024.10.26.620382}{10.1101/2024.10.26.620382}

    \item Sorjonen K, Melin B, Nilsonne G. Still spurious: A comment on attempts to revive cognitive ability tilts. doi:~\href{https://doi.org/10.31234/osf.io/tbhx6}{10.31234/osf.io/tbhx6}

    \item Sorjonen K, Melin B, Nilsonne G. The Spurious Prospective Associations Model (SPAM): Explaining longitudinal associations due to statistical artifacts.

    \item Sorjonen K, Melin B, Nilsonne G. Inconclusive evidence for a prospective effect of academic self-concept on achievement: A simulated reanalysis and comment on Marsh et al. (2024). doi:~\href{https://doi.org/10.31234/osf.io/3cjtr}{10.31234/osf.io/3cjtr}

    \item Sorjonen K, Nilsonne G, Melin B. Triangulation can be useful depending on what you wish to achieve: A reply to Lucas et al. (2024). doi:~\href{https://doi.org/10.31234/osf.io/6dq2j}{10.31234/osf.io/6dq2j}

    \item Nilsonne G. \textit{Cross-lagged network models do not prove causality: Results from simulations and analyses of symptoms of depression and anxiety.}. doi:~\href{https://doi.org/10.31234/osf.io/j3k68}{10.31234/osf.io/j3k68}

\end{enumerate}


%% ================================================================
%% BOOKS
%% ================================================================

\section{Books}
\begin{enumerate}[nolistsep, label={[\arabic*]}]

    \item Affective Regulation and the Neuroscience of Emotion: Festschrift in Honor of Arne Öhman. 2011.

\end{enumerate}


%% ================================================================
%% OPEN PEER REVIEWS
%% ================================================================

\section{Open Peer-Review Reports}
\begin{enumerate}[nolistsep, label={[\arabic*]}]

    \item Smith R and Roberts I. Time for sharing data to become routine: the seven excuses for not doing so are all invalid [version 1; referees: 2 approved, 1 approved with reservations] F1000Research 2016, 5:781. doi: 10.12688/f1000research.8422.1. doi of review report: 10.5256/f1000research.9066.r14294

    \item Matthews PC. Fairness in scientific publishing [version 2; referees: 3 approved] F1000Research 2016, 5:2816. doi: 10.12688/f1000research.10318.2. doi of review report 1: 10.5256/f1000research.11114.r18242. doi of review report 2: 10.5256/f1000research.11479.r19405

    \item Rowhani-Farid A, Barnett AG. Badges for sharing data and code at Biostatistics: an observational study [version 1; referees: 1 approved] F1000Research 2018, 7:90. doi: 10.12688/f1000research.13477.1. doi of review report: 10.5256/f1000research.14635.r30171.

\end{enumerate}


%% ================================================================
%% REPORTS
%% ================================================================

\section{Reports}
\begin{enumerate}[nolistsep, label={[\arabic*]}]

    \item Omskärelse av pojkar och män: risker och potentiella hälsovinster. 2015.

    \item Öppen tillgång till forskningsdata och publikationer – möjligheter och utmaningar. 2018.

    \item Recommendations and Guidelines on data sharing. 2020. doi:~\href{https://doi.org/10.15497/rda00052}{10.15497/rda00052}

    \item Datahanteringsplaner och policy för öppen tillgång till forskningsdata. 2020.

    \item Fostering Open Science in Europe: Engagement Strategies from EOSC's Task Forces on Research Careers and Curricula. 2024. doi:~\href{https://doi.org/10.5281/zenodo.10925393}{10.5281/zenodo.10925393}

\end{enumerate}


%% ================================================================
%% STUDY MATERIALS
%% ================================================================

\section{Study Materials}
\begin{enumerate}[nolistsep, label={[\arabic*]}]

    \item Nilsonne G, Arborelius L. doi:~\href{https://doi.org/10.6084/m9.figshare.5782944}{10.6084/m9.figshare.5782944}

\end{enumerate}


%% ================================================================
%% DIGITAL RESEARCH OBJECTS
%% ================================================================

\section{Other Digital Research Objects}
\begin{enumerate}[nolistsep, label={[\arabic*]}]

    \item Nilsonne G. Interactive classroom memory tests using online tools. doi:~\href{https://doi.org/10.6084/m9.figshare.1449013}{10.6084/m9.figshare.1449013}

    \item Beyer F, Flannery J, Gau R, Janssen L, Schaare L, Hartmann H, et al. A fMRI pre-registration template.

    \item Nilsonne G. Figure: interlinked digital research objects. doi:~\href{https://doi.org/10.53962/hbfx-kmg3}{10.53962/hbfx-kmg3}

    \item Nilsonne G, Ahnström, L. Digitization of data from published plots v1. 0. doi:~\href{https://doi.org/10.17504/protocols.io.n2bvj8wnngk5/v1}{10.17504/protocols.io.n2bvj8wnngk5/v1}

\end{enumerate}


%% ================================================================
%% SCHOLARLY DEBATE
%% ================================================================

\section{Scholarly Debate}
\begin{enumerate}[nolistsep, label={[\arabic*]}]

    \item Nilsonne G. Open access ökar forskningens tillgänglighet. \textit{Läkartidningen}. 2010.

    \item Nilsonne G. Läkarstudenters minskande empati under utbildningen ifrågasätts. \textit{Läkartidningen}. 2010.

    \item Nilsonne G. Om bevisföring i klinisk medicin och i allmänhet. \textit{Folkvett}. 2011.

    \item Nilsonne G. Towards an ecosystem for open data. \textit{BMJ blogs}. 2014.

    \item Nilsonne G. Våga vårda vårt fackspråk. \textit{Läkartidningen}. 2014.

    \item Nilsonne G. Öppna data viktig pusselbit i framtidens forskningsmetod. \textit{Läkartidningen}. 2014.

    \item Nilsonne G. Clinical trials, replication, and crowdsourcing. \textit{Open Science Collaboration blog}. 2014.

    \item Nilsonne G. Fusk bekämpas bäst med öppna forskningsdata. \textit{Curie}. 2015.

    \item Nilsonne G. Mindre än hälften av psykologistudier kunde reproduceras. \textit{Läkartidningen}. 2015.

    \item Nilsonne G. Bias i metaanalys kan leda till överskattade effekter. \textit{Läkartidningen}. 2016.

    \item Nilsonne G. Willén R. Öppna data kan öka kunskapsmassan och motverka fusk. \textit{Läkartidningen}. 2016.

    \item Nilsonne G. Vetenskapsrådet bör agera mer kraftfullt. 2016.

    \item Kanstrup M, Moberg K, Gustafsson S, Gornitzki C, Nilsonne G, Galanti MR. Upprop för välgrundad forskning och vård. \textit{KI-bladet}. 2016.

    \item Willén R, Nilsonne G. Psykologin i framkant för öppen vetenskap. \textit{Psykologtidningen}. 2017.

    \item Dietrichson E, Nilsonne G, Palm A. Ingen motsättning mellan demokrati och evidens. \textit{Skola och Samhälle}. 2017.

    \item Nilsonne G. Öppen tillgång ger författare starkare ställning. \textit{Universitetsläraren}. 2017.

    \item Nilsonne G. Open science: providing skills and inspiring change. \textit{Stockholm University Press blog}. 2017.

    \item Nilsonne G. Replik: Befria forskningen från tidskrifternas tvångströja. \textit{Curie}. 2018.

    \item Appelhoff S, Bates JF, Ghosh S, Keator DB, Kennedy DN, Poldrack R, Poline JB, Steffener J, Nichols BN, Feingold F, Pernet C, Nilsonne G, Maumet C, Flandin G, Gau R, Oostenveld R, Dupre E, Delorme A, Markiewicz CJ, Perez N, Helmer KG, Jarecka D, Grethe JS, Patterson D, Auer T, Bartsch H, Nichols TE, Calhoun V. BIDS and the NeuroImaging Data Model (NIDM). \textit{F1000Research}. 2019.

    \item Nilsonne G. Dyrköpt avtal med Elsevier. \textit{Curie}. 2019.

    \item Nilsonne G. Stor variation i resultat när olika forskare analyserade samma data. \textit{Läkartidningen}. 2020.

    \item Nilsonne G. Rekommendationer om Covid-19-data från Research Data Alliance (RDA). \textit{SND nyhetsbrev}. 2020.

    \item Brembs B, Nilsonne G, Susi T. How to reclaim ownership of scholarly publishing. 2022.

    \item Brembs B, Huneman P, Schönbrodt F, Nilsonne G, Susi T, Siems R, Perakakis P, Trachana V, Ma L, Rodriguez-Cuadrado S. Creating a market to replace publisher monopolies. \textit{sOApbox}. 2022.

    \item Widmark W, Anderberg S, Holmgren S, Lindvall J, Nilsonne G. Sverige får inte missa tåget när EU satsar på öppen vetenskap. \textit{Curie}. 2022.

    \item Widmark W, Anderberg S, Holmgren S, Lindvall J, Nilsonne G. Forskningsbibliotek viktiga för Europasatsning på öppna data. \textit{Biblioteksbladet}. 2022.

    \item Nilsonne G. Dags att överge traditionella tidskrifter. 2022.

    \item Nilsonne G, Jönsson JI, Söderbergh Widding A. Dags att förnya hur forskning bedöms. \textit{Curie}. 2022.

    \item Nilsonne G. Bakvägsidentifiering av forskningsdata med personuppgifter: riskbedömning och riskminskande åtgärder. \textit{Svepet: medlemsblad för Svensk epidemiologisk förening (SVEP)}. 2022.

    \item Nilsonne G. Replikeringskrisen hotar vetenskapens trovärdighet. \textit{Folkvett}. 2023.

    \item Axfors C, Nilsonne G. Ofullständig rapportering av studier försämrar kunskapsläget. \textit{Läkartidningen}. 2023.

    \item Open Science Community Sverige. 2023.

    \item Nilsonne G, Axfors C, Zavalis E. Utmaningar för den forskande läkaren. \textit{Läkartidningen}. 2023.

    \item Nilsonne G,. ”Vetenskapliga tidskrifter har spelat ut sin roll”. 2023.

    \item Nilsonne G,. Uppföljning av öppen vetenskap – hur gör vi?. \textit{SND Nätverks-nytt}. 2023.

    \item Nilsonne G,. Riktlinjer för öppen vetenskap en besvikelse. \textit{Upsala Nya Tidning}. 2023.

    \item Nilsonne G, O'Neill G, Dahle S, Gaillard V, Priess-Buchheit J, Schmidt B, Psomopoulos F. EOSC as an Enabler of Research Assessment Reform: Position Paper from Task Force on Research Careers, Recognition, and Credit. 2023. [\href{https://doi.org/10.5281/zenodo.10417069}{data}]

    \item Tikkanen KW, Nilsonne G. Delning av forskningsdata – att göra humanistisk forskning nåbar för fler. 2024.

    \item Bishop D, et al. An open letter regarding Scientific Reports. 2024.

\end{enumerate}


%% ================================================================
%% INVITED TALKS
%% ================================================================

\section{Invited Talks}
\begin{enumerate}[nolistsep, label={[\arabic*]}]

    \item Open Access från en forskares horisont. 2011.

    \item Varför open access - en forskares perspektiv. 2011.

    \item Emotional Regulation and Empathy – Pharmacological and Translational Studies. 2012.

    \item Psychopathic traits and the representation of others’ emotional states. 2012.

    \item Partial sleep deprivation and emotional function: the Stockholm Sleepy Brain Project. 2013.

    \item Effects of sleep deprivation on brain functions in the emotional domain. 2014.

    \item Prosocial enhancement. 2014.

    \item Sleep deprivation and the brain: focus on emotion. 2014.

    \item Why do psychopathic traits persist? - an evolutionary perspective. 2015.

    \item Reproducibility in psychological science. 2015.

    \item Psychopathy: Brain mechanisms from an evolutionary perspective. 2015.

    \item Replication science and other methods for reproducible research. 2016.

    \item What’s new in open science and reproducible research?. 2016.

    \item Öppen och reproducerbar vetenskap. 2016.

    \item Open data. 2016.

    \item Open science and reproducible research. 2016.

    \item Open data: examples and implications. 2016.

    \item Öppna och dolda data – konsekvenser för forskning och hälsa. 2016. [\href{https://www.youtube.com/watch?v=2noejs93mY8}{video}]

    \item Open data. 2016.

    \item Öppna data. 2016. [\href{https://www.youtube.com/watch?v=1azYwpiAunA}{video}]

    \item Forskningsetik och öppna data. 2017.

    \item Effects of sleep deprivation on resting state connectivity and emotional processing: results from the Stockholm Sleepy Brain study. 2017.

    \item Brain imaging of sleep deprivation in healthy older humans – results from the Stockholm Sleepy Brain study. 2017.

    \item Reproducibility and Openness in Psychiatry and Psychology, Karolinska Institutet, Centre for Psychiatry Research, 2017-11-17. 2017. [\href{https://osf.io/6zu3s/}{slides}]

    \item Reproducibility and Openness in Psychiatry and Psychology, Karolinska Institutet, Kompetenscentrum för psykoterapi, 2017-12-05. 2017.

    \item Öppen vetenskap – principer, exempel och erfarenheter. 2017.

    \item Open data - examples and reflections. 2018.

    \item Open data in neuroimaging research. 2018. [\href{https://osf.io/s8bep/}{slides}]

    \item Öppna forskningsdata. 2018.

    \item Åtgärder för öppna data. 2018.

    \item Kan vi skilja mellan bra och dålig forskning?. 2018.

    \item Intrinsic brain connectivity after partial sleep deprivation in young and older adults: results from the Stockholm Sleepy Brain study. 2018.

    \item Preprints – a discussion. 2018. [\href{https://osf.io/wfctz/}{slides}]

    \item Reproducibility and Openness in Psychological Science. 2018. [\href{https://osf.io/yre8g/}{slides}]

    \item Open and FAIR data from human participants in neuroscientific research. 2018. [\href{https://osf.io/8jbqh/}{slides}]

    \item Meta-science in neuroscience: examples and perspectives. 2018. [\href{https://osf.io/f6hjs/}{slides}]

    \item Open methodology and reproducible research. 2019. [\href{https://osf.io/ah6sy/}{slides}]

    \item Panel discussion about open science. 2019.

    \item Open Science: a researcher’s perspective. 2019.

    \item Open science – what a young researcher needs to know?. 2019. [\href{https://osf.io/xbcde/}{slides}]

    \item Open neuroimaging data and personal data privacy: convergence or divergence. 2020. [\href{https://www.youtube.com/watch?v=A3BapxKPnRQ&list=PLVso6Qs8PLChTc8uMUWNSe_Bj4tfOHC_f&index=27}{video}]

    \item Pre-registration, Registered reports. 2020. [\href{https://www.youtube.com/watch?v=-GKKiwMyfH8&list=PLVso6Qs8PLChTc8uMUWNSe_Bj4tfOHC_f&index=15}{video}]

    \item Brain research data sharing and personal data privacy: continued discussion on anonymization + community meeting. 2020. [\href{https://www.youtube.com/watch?v=j9wFwBxwMvI&list=PLVso6Qs8PLChTc8uMUWNSe_Bj4tfOHC_f&index=16}{video}]

    \item Publicly sharing data – opportunities and challenges. 2020. [\href{https://osf.io/tp37s/}{slides}]

    \item Öppen data och registerforskning. 2020. [\href{https://osf.io/sfn4v/}{slides}]

    \item Open and FAIR data – some principles and examples. 2021. [\href{https://osf.io/3at4s/}{slides}]

    \item Betydelsen av datahanteringsplaner. 2021. [\href{https://osf.io/nmdp4/}{slides}]

    \item Is PrePublication the future? Challenges and opportunities after a pandemic. 2021.

    \item Reproducibility in science –challenges and solutions. 2021. [\href{https://osf.io/9uw4x/}{slides}]

    \item The multi-analyst design - examples and outstanding metascientific questions. 2021. [\href{https://osf.io/yprkw/}{slides}]

    \item Replacing Academic journals. 2021. [\href{https://osf.io/uktqb/}{slides}]

    \item "It matters how we open knowledge": A panel discussion on open access and equity. 2021. [\href{https://www.youtube.com/watch?v=6oZbVTqhX5o}{video}]

    \item Methods to assess and mitigate re-identification risks when sharing research data. 2021. [\href{https://osf.io/7xenk/}{slides}]

    \item How to reclaim ownership of scholarly publishing. 2022. [\href{https://osf.io/rc82z/}{slides}]

    \item Meritvärdering för ett öppet forskningssystem i Europa. 2022. [\href{https://osf.io/dqyp3/}{slides}]

    \item Big team science in Psychology - How to use large samples and many analysts to assess and improve reproducibility. 2022. [\href{https://osf.io/b3wh6/}{slides}]

    \item Using the Brain Imaging Data Structure (BIDS) to organise and analyse data. 2022. [\href{https://osf.io/4rtcx/}{slides}]

    \item Reproducerbarhetskrisen i biomedicinsk forskning. 2022.

    \item Implementering av SUHF:s färdplan för öppen vetenskap. 2022. [\href{https://osf.io/xf9nk/}{slides}]

    \item Replacing Academic Journals. 2022. [\href{https://osf.io/amz5w/}{slides}]

    \item Research Reproducibility in the Social Sciences. 2022. [\href{https://osf.io/zvm7p/}{slides}]

    \item I stället för tidskrifter. 2022.

    \item Omställning till ett öppet vetenskapssystem: svenska lärosätens roll. 2022. [\href{https://osf.io/zvyqa/}{slides}]

    \item Open Science. 2022.

    \item Öppen tillgång till forskningsdata. 2022. [\href{https://osf.io/s26xh}{slides}]

    \item Pathways towards community-based, transparent, and systematic quality control of research objects. 2022. [\href{https://osf.io/4u6jc}{slides} | \href{https://www.youtube.com/watch?v=QqGHW5tb87s}{video}]

    \item Perspectives on scientific reproducibility. 2022. [\href{https://osf.io/hnju6}{slides}]

    \item Transparency and reproducibility in neuroscience. 2022. [\href{https://osf.io/nx2yw}{slides}]

    \item Transparency and reproducibility in psychiatry research. 2022. [\href{https://osf.io/wnm4r}{slides}]

    \item Systematic, transparent, and community-based assessment and quality control of research data. 2022. [\href{https://osf.io/wnm4r}{slides}]

    \item Open science and data sharing. 2022. [\href{https://osf.io/bvu2s}{slides}]

    \item Teaching open science and reproducible research at Karolinska Institutet. 2022. [\href{https://osf.io/9s2pd}{slides}]

    \item How to make science more trustworthy by improving transparency and reproducibility. 2023. [\href{https://osf.io/ung8q}{slides}]

    \item Historien bakom Coara. 2023. [\href{https://osf.io/e8pvs}{slides}]

    \item The why and how of data sharing. 2023. [\href{https://osf.io/ht7d5}{slides}]

    \item Bortom transformativa avtal. 2023. [\href{https://osf.io/85a6x}{slides}]

    \item Open science and reproducible research. 2023. [\href{https://osf.io/5a8f2}{slides}]

    \item Brainstorming new ways to drive clinical trial reporting. 2023. [\href{https://www.youtube.com/watch?v=HohWoE3bMi0}{video}]

    \item Debate: Should sponsors be sanctioned for failing to report clinical results?. 2023. [\href{https://www.youtube.com/watch?v=IjsGeVlkXeE}{video}]

    \item Teaching reproducible research practices – experiences from Karolinska Institutet. 2023. [\href{https://osf.io/e6vrj}{slides}]

    \item Sekundäranvändning av forskningsdata – forskarens perspektiv. 2023. [\href{https://osf.io/xthbf}{slides}]

    \item CoARA – bakgrund och innebörd ur en forskares perspektiv. 2023. [\href{https://osf.io/85a6x}{slides}]

    \item Rapportering av utförda studier – en aspekt av forskningens transparens och integritet. 2023. [\href{https://osf.io/jc7dy}{slides}]

    \item Data sharing in neuroimaging research. 2023. [\href{https://osf.io/fk3ny}{slides}]

    \item Open and accessible data – a recipe for research success?. 2023.

    \item Ett forskarperspektiv. 2023.

    \item Clinical trials reporting at Nordic medical universities and university hospitals. 2023. [\href{https://www.youtube.com/watch?v=RoYXvmG0KN0}{video} | \href{https://osf.io/z2cyn}{slides}]

    \item Pathways to an open science system: Replacing academic journals. 2023. [\href{https://osf.io/jvw5g}{slides}]

    \item Replacing Academic Journals. 2023.

    \item Replacing Academic Journals. 2023.

    \item Clinical trials reporting at Nordic medical universities and university hospitals. 2023. [\href{https://osf.io/bh98f}{slides} | \href{https://www.youtube.com/watch?v=GBNQbQvYzjc}{video}]

    \item Clinical trials reporting at Nordic medical universities and university hospitals. 2023. [\href{https://osf.io/cdqm6}{slides}]

    \item ReproducibiliTea Journal Club. 2024.

    \item On Metascience. 2024. [\href{https://osf.io/xq623}{slides}]

    \item Metascience in medical research–some examples. 2024. [\href{https://osf.io/2xq5h}{slides}]

    \item Så kan medicinsk forskning bli mer pålitlig och användbar. 2024. [\href{https://osf.io/j7tgk}{slides}]

    \item Research Assessment for Transparency and Reproducibility. 2024.

    \item Multi-analyst research designs and the influence of analytical variability on the robustness of research results. 2024. [\href{https://osf.io/rbshd}{slides}]

    \item Data sharing in research on human participants. 2024. [\href{https://osf.io/79hcn}{slides}]

    \item The long road towards open science: strategies for change. 2024. [\href{https://osf.io/rn8wf}{slides} | \href{https://www.youtube.com/watch?v=svJvq1cot5k}{video}]

    \item On Metascience. 2024.

    \item Metascience for better research transparency and reproducibility. 2024. [\href{https://osf.io/9wq8c}{slides}]

    \item Data sharing – some principles and examples. 2024. [\href{https://osf.io/cw5bn}{slides}]

    \item Varför öppna forskningsdata? Juridiken som möjliggörare. 2024. [\href{https://osf.io/ypwm5}{slides}]

    \item Concepts in Open Science. 2025. [\href{https://osf.io/5r8v9}{slides}]

    \item Öppna forskningsdata för ett rationellt samhälle. 2025.

\end{enumerate}


%% ================================================================
%% CONFERENCE PRESENTATIONS
%% ================================================================

\section{Conference Presentations}
\begin{enumerate}[nolistsep, label={[\arabic*]}]

    \item Nilsonne G. Title not on record. Läkaresällskapets Riksstämma, Sektionen för Patologi, Stockholm 2003. Oral presentation. 2003.

    \item Nilsonne G, Sun X, Nyström C, Rundlöf A-K, Björnstedt M, Dobra K. Selenite induces apoptosis in malignant mesothelioma cells. 2004.

    \item Nilsonne G, Sun X, Rundlöf AK, Fernandes AP, Nyström C, Stein A, Kocic B, Björnstedt M, Hjerpe A, Dobra K. Selenite induces apoptosis in therapy resistant malignant mesothelioma cells. 2005.

    \item Nilsonne G, Sun X, Rundlöf AK, Fernandes AP, Nyström C, Stein A, Kocic B, Björnstedt M, Hjerpe A, Dobra K. Selenite induces apoptosis in therapy resistant malignant mesothelioma cells. 2005.

    \item Nilsonne G et al. Title not on record. Cell Signalling World: Signal Transduction Pathways as Therapeutic Targets, Luxembourg 2006. Poster presentation. 2006.

    \item Nilsonne G et al. Title not on record. KI Cancer Retreat, Djurönäset 2006. Poster presentation. 2006.

    \item Nilsonne G, Stein A, Rundlöf AK, Fernandes AP, Björnstedt M, Hjerpe A, Dobra K. Targeting thioredoxin reductase - a possible new treatment for malignant mesothelioma. 2006.

    \item Nilsonne G et al. Title not on record. KI Cancer Retreat, Djurönäset 2007. Poster presentation. 2007.

    \item Nilsonne G, Olm E, Szulkin A, Mundt F, Stein A, Kosic B, Rundlöf AK, Fernandes AP, Björnstedt M, Dobra K. Selenite-induced apoptosis signalling in differentially sensitive mesothelioma cell lines. 2007.

    \item Nilsonne G et al. Title not on record. KI Cancer Retreat, Djurönäset 2008. Poster presentation. 2008.

    \item Nilsonne G et al. Title not on record. KI Cancer Retreat, Djurönäset 2009. Poster presentation. 2009.

    \item Nilsonne G, Tamm S, Golkar A, Gospic K, Olsson A, Ingvar M, Petrovic P. Effects of 25 mg Oxazepam on self-reported and psychophysiological measures of empathy. 2012.

    \item Nilsonne G, Tamm S, Golkar A, Gospic K, Olsson A, Ingvar M, Petrovic P. Effects of 25 mg Oxazepam on self-reported and psychophysiological measures of empathy. 2012.

    \item Nilsonne G, Tamm S, Golkar A, Gospic K, Olsson A, Ingvar M, Petrovic P. Effects of 25 mg Oxazepam on empathy: a behavioural study. 2012.

    \item Nilsonne G, Sörman K, Tamm S, Golkar A, Gospic K, Olsson A, Kristiansson M, Ingvar M, Petrovic P. Psychopathic traits and the representation of others’ emotional states. 2013.

    \item Nilsonne G, Tamm S, Golkar A, Gospic K, Olsson A, Ingvar M, Petrovic P. Pharmacological modulation of empathy using Oxazepam. 2013.

    \item Nilsonne G, Tamm S, d’Onofrio P, Schwarz J, Kecklund G, Åkerstedt T, Fischer H, Lekander M. Effect of partial sleep deprivation on self-rated health and sickness. 2013.

    \item Nilsonne G, Tamm S, d’Onofrio P, Schwarz J, Kecklund G, Lekander M, Åkerstedt T, Fischer H. Detection of facial mimicry by electromyography during fMRI scanning. 2013.

    \item Nilsonne G, Tamm S, d’Onofrio P, Thuné H, Schwarz J, Petrovic P, Fischer H, Kecklund G, Åkerstedt T, and Lekander M. Effect of partial sleep deprivation on empathic responding in an fMRI experiment. 2014.

    \item Nilsonne G, Tamm S, d’Onofrio P, Thuné H, Schwarz J, Petrovic P, Fischer H, Kecklund G, Åkerstedt T, and Lekander M. Effect of partial sleep deprivation on empathy for pain in an fMRI experiment. 2014.

    \item Nilsonne G. Sleep Deprivation and Emotion – fMRI studies. 2014.

    \item Nilsonne G. Effects of Sleep Deprivation on Emotional Processing Related to Others’ Emotional Expression and Experience. 2015.

    \item Tamm S, et al. It hurts me too: an fMRI study on the effects of experimental sleep restriction on empathy for pain in younger and older adults. 2016.

    \item Nilsonne G. Öppna, slutna och dolda data: erfarenheter av en metaanalys och synpunkter från en forskares horisont. 2016. [\href{http://www.slideshare.net/kbredaktion/ppna-slutna-och-dolda-data-erfarenheter-av-en-metaanalys-och-synpunkter-frn-en-forskares-horisont-gustav-nilsonne}{slides} | \href{https://www.youtube.com/watch?v=jEfRBOV_KEA\&index=14\&list=PLZVkEICvA5-HRSSr5kb-hswTHUVag5c8P}{video}]

    \item Nilsonne G, commentator on presentation by Gravert C:. The hidden costs of nudging: Experimental evidence from reminders in fundraising. 2016.

    \item Nilsonne G. Panelsamtal. Öppen Vetenskap – Möjligheter och utmaningar, Stockholm 2017. Program: \href\{http://www.suhf.se/storage/ma/d3809593a7aa480e974101d18a93cbe9/15fb3e16f5. 2017.

    \item Nilsonne G, Tamm S, Schwarz J, Almeida R, Fischer H, Kecklund G, Lekander M, Fransson P, Åkerstedt T. Increased global fMRI signal variability after partial sleep deprivation: findings from the Stockholm Sleepy Brain study. 2017.

    \item Nilsonne G, Inge M, Hedlund M, Axelsson J, Åkerstedt T, Lekander M. Interleukin-6: diurnal variation and relation to sleep. 2017. [\href{https://osf.io/kd53f/}{slides}]

    \item Nilsonne G, Tamm S, Schwarz J, Almeida R, Fischer H, Kecklund G, Lekander M, Fransson P, Åkerstedt T. Increased global fMRI signal variability after partial sleep deprivation: findings from the Stockholm Sleepy Brain study. 2017.

    \item Nilsonne G, Gilmore R. Workshop: IRB for open data. 2017.

    \item Nilsonne G. Salvaging data: retrospective publication. 2017.

    \item Nilsonne G. Incentivizing open data with badges. 2017. [\href{https://osf.io/r7xv2/}{slides} | \href{https://www.youtube.com/watch?v=FztK2AG0W4Y&feature=youtu.be}{video}]

    \item Nilsonne G. The FAIR principles and global health research. 2018. [\href{https://osf.io/2nxmr/}{slides}]

    \item Nilsonne G. Dissipating and cumulating research fields - new concepts for promoting data salvage. 2018. [\href{https://osf.io/u7z23/}{slides}]

    \item Nilsonne G. Reproducerbarhet i psykologisk och psykiatrisk forskning - några exempel. 2019. [\href{https://osf.io/5nt6s/}{slides}]

    \item Nilsonne G. Reproducerbarhet i psykologisk och psykiatrisk forskning - några exempel. 2019. [\href{https://osf.io/n6mhf/}{slides}]

    \item Nilsonne G. Open data in research on humans – promises and challenges. 2019. [\href{https://osf.io/ja4sr/}{slides}]

    \item Nilsonne G. Anonymization and secondary use of individual participant data. 2020.

    \item Nilsonne G. Data sharing for scientific integrity - examples and frameworks. 2020. [\href{https://osf.io/sz3wn/}{slides}]

    \item Nilsonne G. Ethics of FAIR in Neuroimaging. 2021. [\href{https://osf.io/bswxp/ }{slides}]

    \item Nilsonne G. Data sharing for clinical research – challenges and opportunities. 2021. [\href{https://osf.io/4rgj3/}{slides}]

    \item Nilsonne G. Data domain specialists and flagship projects: challenges and successes in crafting partner contributions to the Swedish National Data Service consortium. 2022.

    \item Nilsonne G. Changing the publication system. 2022.

    \item Nilsonne G. Olika slutsatser från samma data - analytisk flexibilitet i medicinsk forskning. 2023. [\href{https://osf.io/fm7we}{slides}]

    \item Nilsonne G. EEGManyPipelines: A large-scale, grass-root multi-analyst study of EEG analysis practices. 2023. [\href{https://osf.io/zn2vf}{slides}]

    \item Natalie Hyltse, Rickard Carlsson, Nilsonne G. Risk Of Significant Error - A checklist to assess the risk of errors in statistical reporting and datasets. 2023.

    \item Nilsonne G. Ensuring DeSci results are reproducible and trustworthy. 2023. [\href{https://osf.io/z5amvh}{slides}]

    \item Nilsonne G. Research assessment as a lever to improve scientific practices. 2023. [\href{https://vimeo.com/showcase/10628629/video/862484008}{video}]

    \item Nilsonne G. Clinical trials reporting at Nordic medical universities and university hospitals. 2023. [\href{https://osf.io/z2cyn}{slides}]

    \item Nilsonne G. Coara – vad är det?. 2023. [\href{https://osf.io/2jw5n}{slides}]

    \item Nilsonne G. Big team science reproducibility collaborations and transnational initiatives. 2023. [\href{https://osf.io/b7v3y}{slides}]

    \item Natalie Hyltse, Rickard Carlsson, Nilsonne G. Risk Of Statistical Error - A checklist to assess the risk of errors in statistical reporting and datasets. 2023.

\end{enumerate}


%% ================================================================
%% CONFERENCE ABSTRACTS
%% ================================================================

\section{Other Indexed Conference Abstracts}
\begin{enumerate}[nolistsep, label={[\arabic*]}]

    \item Nilsonne G, Sun X, Nyström C, Rundlöf A-K, Nurminen M, Björnstedt M, Dobra K. Differential effects of selenite on malignant mesothelioma cells - a novel therapeutic approach related to the Trx-TrxR system and the redox state of the cells. 2004.

    \item Tamm S, Nilsonne G, d'Onofrio P, Thuné H, Schwarz J, Petrovic P, Fischer H, Kecklund G, Åkerstedt T, Lekander M. Effect of partial sleep deprivation on empathy for pain in an fMRI experiment. 2014.

    \item Thuné H, Nilsonne G, Tamm S, Schwarz J, Kecklund G, Lekander M, Åkerstedt T, Fischer, H. Effects of partial sleep deprivation on the neural mechanisms of face perception. 2014.

    \item Åkerstedt T, Lekander M, Nilsonne G, Fischer H, Kecklund G, d'Onofrio P, Gruber G, Schwarz J. Women sleep better and have a stronger response to late night curtailed sleep than men, particularly in older individuals-effects on polysomnographical sleep. 2016.

    \item Månsson K, Lindqvist D, Yang L, Wolkowitz O, Nilsonne G, Isung J, Svanborg C, Boraxbekk C-J, Fischer H, Lavebratt C, Furmark T. S13. Can Psychological Treatment Slow Down Cellular Aging in Social Anxiety Disorder? An Intervention Study Evaluating Changes in Telomere Length and Telomerase Activity. 2018.

    \item Åkerstedt T, Lekander M, Nilsonne G, Tamm S, d’Onofrio P, Kecklund G, Fischer H, Petrovic P, Månsson KN. 0149 Gray Matter Volume Correlates Of Sleepiness: A Voxel-based Morphometry Study In Younger And Older Adults. 149.

    \item Nilsonne G, Tamm S, Schwarz J, Kecklund G, Lekander M, Petrovic P, Åkerstedt T. Sleepy Brain 2: Multimodal brain imaging after a night of sleep deprivation. 2018. [\href{https://osf.io/vxmdj/}{slides}]

\end{enumerate}


%% ================================================================
%% POPULAR SCIENCE
%% ================================================================
\newpage
\section*{Popular Science}


\section{Blogging}
\begin{enumerate}[nolistsep, label={[\arabic*]}]

    \item Evolving ideas, personal blog, 2009. url: evolvingideas.wordpress.com/

    \item nilsonne.net, 2017-2018.

    \item Risk of bias in a meta-analysis of cytokine-inhibiting drugs against depression: reflections from the 24th PNIRS meeting. archived: figshare, doi: 10.6084/m9.figshare.5027387.v1

    \item A tentative assessment of fMRI data sharing in the journal Plos One: results from the Stockholm BrainHack, with Stefan Wiens. archived: figshare, doi: 10.6084/m9.figshare.5027384.v1

\end{enumerate}



\section{Popular Science Talks}
\begin{enumerate}[nolistsep, label={[\arabic*]}]

    \item Sex och kärlek i hjärnan. 2011.

    \item Varför behöver vi en ny sexualvaneundersökning?. 2011.

    \item Sex och kärlek i hjärnan. 2011.

    \item Kärlek och sex i hjärnan. 2012.

    \item Annat sex. 2012.

    \item Empati - när, var, och hur?. 2012.

    \item Informationsinfrastruktur för ett rationellt samhälle. 2012.

    \item Sömn, sömnbrist och emotionell reglering. 2014. [\href{https://www.youtube.com/watch?v=wq2C_KNrn-8}{video}]

    \item How can we live forever? The dilemma of human enhancement. 2015. [\href{https://www.youtube.com/watch?v=74irIhJVKK8\&index=32\&list=PLvm-fAk-DoHVR6aWtbjW3oX9J7\_3KIhH3}{video}]

    \item Evidensbaserad politik. 2015.

    \item Evidensbaserad policy. 2015.

    \item Vi måste randomisera oss. 2015.

    \item Oxazepam and emotional responding. 2015.

    \item Stockholm Pride/Vetenskap och Folkbildning 2015-07-27: Häng med en forskare. 2015.

    \item Radio bubb.la.

    \item Vetenskapsspecial 2015-02-14. SoundCloud: radiobubbla. 2015.

    \item Vetenskapsspecial 2016-10-20. SoundCloud: radiobubbla. 2016.

    \item Vetenskapsspecial 2016-11-18. SoundCloud: radiobubbla. 2016.

    \item Medicinspecial 2016-12-13. SoundCloud: radiobubbla. 2016.

    \item Vetenskapsspecial 2016-12-28. SoundCloud: radiobubbla. 2016.

    \item Panelsamtal Människa Plus 2016-11-27. Web: kulturhusetstadsteatern.se. 2016.

    \item Panelsamtal tidningen Curies femårsjubileum 2016-11-30. 2016.

    \item Red Vic lecture: The Reproducibility Crisis 2017-08-15. Web: facebook.com. 2017.

    \item Panelsamtal Axfoundation: Hotet mot sanningen 2017-10-10. 2017.

    \item Replikeringsproblemet och dess möjliga lösningar. 2018.

    \item Replikeringsproblemet och dess möjliga lösningar. 2018.

    \item Ett forskningsprojekt från ax till limpa. 2019. [\href{https://osf.io/zx59b/}{slides}]

    \item Demokratisera forskningen - till varje pris?. 2019. [\href{https://www.facebook.com/vetenskapoallm/videos/882241585458235/}{video}]

    \item A research project from start to finish. 2019.

    \item Forskning och verklighet – vilka vetenskapliga påståenden kan vi lita på och varför?. 2019. [\href{https://osf.io/y65gr/}{slides}]

    \item Gustav Nilsonne om replikerbarhet och open science. 2021.

    \item Forskning vi kan lita på - en förutsättning för rationellt beslutsfattande. 2022. [\href{https://osf.io/cbtgq/}{slides}]

    \item Falska forskningsfynd?. 2022.

    \item Öppen vetenskap och forskningens pålitlighet. 2022.

    \item Meta-vetenskap med Gustav Nilsonne. 2023.

    \item Vägen mot öppen tillgång till forskningsdata - hur går omställningen?. 2024. [\href{https://www.youtube.com/watch?v=pa9PV9jw5-M}{video}]

\end{enumerate}



\section{Popular Science Writings and General Debate}
\begin{enumerate}[nolistsep, label={[\arabic*]}]

    \item .

    \item Förförande diagnoser. 2010. [\href{http://www.ottar.se/artiklar/f-rf-rande-diagnoser}{web}]

    \item Ordförandekrönikor i RFSU Stockholms nyhetsbrev 2011-2012:. 2011.

    \item Tankar kring det nya verksamhetsåret. 2011. [\href{http://www.rfsu.se/sv/RFSU-nara-dig/RFSU-Stockholm/Nyhetsbrev/Mars-2011/Ordforandens-rader/}{web}]

    \item Var är alla flickorna?. [\href{http://www.rfsu.se/sv/RFSU-nara-dig/RFSU-Stockholm/Nyhetsbrev/Nummer-3/Var-ar-alla-flickorna/}{web}]

    \item Stora steg mot ökat medlemsinflytande på årets RFSU-kongress. 2011. [\href{http://www.rfsu.se/sv/RFSU-nara-dig/RFSU-Stockholm/Nyhetsbrev/Nummer-4-2011/Stora-steg-mot-okat-medlemsinflytande-pa-arets-RFSU-kongress/}{web}]

    \item Pridefestivalen en höjdpunkt på året. 2011. [\href{http://www.rfsu.se/sv/RFSU-nara-dig/RFSU-Stockholm/Nyhetsbrev/Nummer-5-2011/Pridefestivalen-en-hojdpunkt-pa-aret/}{web}]

    \item Dags för handling i tvångssteriliseringsfrågan. 5201. [\href{http://www.rfsu.se/sv/RFSU-nara-dig/RFSU-Stockholm/Nyhetsbrev/Nyhetsbrev-52011/Dags-for-handling-i-tvangssteriliseringsfragan/}{web}]

    \item Bättre kunskaper krävs om ungas sexuella hälsa. 7201. [\href{http://www.rfsu.se/sv/RFSU-nara-dig/RFSU-Stockholm/Nyhetsbrev/Nyhetsbrev-72011/Battre-kunskaper-kravs-om-ungdomars-sexuella-halsa/}{web}]

    \item Kunskap – en kungsväg till förändring. 7201. [\href{http://www.rfsu.se/sv/RFSU-nara-dig/RFSU-Stockholm/Nyhetsbrev/Nyhetsbrev-72011/Battre-kunskaper-kravs-om-ungdomars-sexuella-halsa/}{web}]

    \item Dags för regeringen att sätta ned foten om tvångssteriliseringarna. 1201. [\href{http://www.rfsu.se/sv/RFSU-nara-dig/RFSU-Stockholm/Nyhetsbrev/Nyhetsbrev-12012/Dags-for-regeringen-att-satta-ned-foten-om-tvangssteriliseringarna/}{web}]

    \item Det ideella engagemanget är öppenhetens grundval. 2011. [\href{https://drive.google.com/open?id=0B79wnU8V_RoiVHRELWtPc0tSdGs}{web}]

    \item Transsexuella har legitim rätt att vilja bli föräldrar. 2012. [\href{https://drive.google.com/open?id=0B79wnU8V_Roic1QySHZnRHJqY3c}{web}]

    \item Studie stärker argument mot omskärelse av pojkar. 2012. [\href{http://www.ottar.se/artiklar/studier-st-rker-argument-mot-omsk-relse-av-pojkar}{web}]

    \item Pride är en angelägenhet för hela Stockholm. 2012. [\href{http://www.dn.se/debatt/pride-ar-en-angelagenhet-for-hela-stockholm/}{web}]

    \item Den fantastiska hjärnan. [\href{http://fjardeuppgiften.se/tag/gustav-nilsonne/}{web}]

    \item Kärlekens kemi: Kärlek blir till i hjärnan. 2012.

    \item Biologiska grunder för sexuell läggning. 2013. [\href{http://www.rfsu.se/Bildbank/Dokument/Metod-Handledning/Den\%20onaturliga\%20naturen_MQ2.pdf?epslanguage=sv}{web}]

    \item Krönikor/gästblogginlägg i Curie 2014-2015:. 2014.

    \item Innovationsfokus eller fri grundforskning – hur blir forskningspolitiken efter valet?. [\href{https://www.tidningencurie.se/kronikor/innovationsfokus-eller-fri-grundforskning---hur-blir-forskningspolitiken-efter-valet}{web}]

    \item Öppna data – framtidens väg. [\href{https://www.tidningencurie.se/kronikor/oppna-data---framtidens-vag}{web}]

    \item Are You an Illusion – betraktelser över svunna kontroverser. [\href{https://www.tidningencurie.se/kronikor/are-you-an-illusion---betraktelser-over-svunna-kontroverser}{web}]

    \item Kan en biolog laga en radio?. [\href{https://www.tidningencurie.se/kronikor/kan-en-biolog-laga-en-radio}{web}]

    \item Bra förslag om open access men kunskapsvinsten måste vara tydlig. [\href{https://www.tidningencurie.se/kronikor/bra-forslag-om-open-access-men-kunskapsvinsten-maste-vara-tydlig}{web}]

    \item Cellteorin och medvetandets gåta. [\href{https://www.tidningencurie.se/kronikor/cellteorin-och-medvetandets-gata}{web}]

    \item Sprutbytesprogram är en fråga om rättvisa. 2015. [\href{http://www.ottar.se/artiklar/sprutbytesprogram-r-en-fr-ga-om-r-ttvisa}{web}]

    \item Creativity, art and science. 2015. [\href{hhttp://vision-forum.blogspot.se/2015/02/creativity-art-and-science.html}{web}]

    \item Kunskap måste gå före ideologi och populism. 2015. [\href{http://www.dn.se/debatt/kunskap-maste-ga-fore-ideologi-och-populism/}{web}]

    \item Evidensbaserad policy. 2015. [\href{https://drive.google.com/file/d/0B6Ii-rt5VieWajZRQjVEdDVHWnM/edit}{web}]

    \item Rösträttsreform måste testas på rätt sätt. 2016. [\href{http://www.dagenssamhalle.se/debatt/roestraettsreform-maste-testas-pa-raett-saett-24121}{web}]

    \item Evidensbaserad policy ger ideologin tydligare roll. 2016. [\href{http://www.frihetssmedjan.se/2016/04/gustav-nilsonne-med-dr-evidensbaserad-policy-ger-ideologin-tydligare-roll/}{web}]

    \item Evidensbasering är ett ungt och innehållsrikt begrepp. 2016. [\href{http://www.frihetssmedjan.se/2016/05/replik-evidensbasering-ar-ett-ungt-och-innehallsrikt-begrepp/}{web}]

    \item Skolor, delta inte i ett dåligt betygsexperiment. 2016. [\href{http://www.svd.se/skolor-delta-inte-i-ett-daligt-experiment}{web}]

    \item Rothstein överdriver problem med evidensbaserad policy. 2019. [\href{http://evidensbaseradpolicy.se/2019/11/12/rothstein-overdriver-problem-med-randomiserade-provningar-av-policy/}{web}]

    \item De utmattade måste få tillgång till rätt vård. 2022.

    \item ”Många forskare struntar i att redovisa sina resultat”. 2024.

\end{enumerate}


\end{document}